% Options for packages loaded elsewhere
\PassOptionsToPackage{unicode}{hyperref}
\PassOptionsToPackage{hyphens}{url}
%
\documentclass[
]{article}
\usepackage{amsmath,amssymb}
\usepackage{iftex}
\ifPDFTeX
  \usepackage[T1]{fontenc}
  \usepackage[utf8]{inputenc}
  \usepackage{textcomp} % provide euro and other symbols
\else % if luatex or xetex
  \usepackage{unicode-math} % this also loads fontspec
  \defaultfontfeatures{Scale=MatchLowercase}
  \defaultfontfeatures[\rmfamily]{Ligatures=TeX,Scale=1}
\fi
\usepackage{lmodern}
\ifPDFTeX\else
  % xetex/luatex font selection
\fi
% Use upquote if available, for straight quotes in verbatim environments
\IfFileExists{upquote.sty}{\usepackage{upquote}}{}
\IfFileExists{microtype.sty}{% use microtype if available
  \usepackage[]{microtype}
  \UseMicrotypeSet[protrusion]{basicmath} % disable protrusion for tt fonts
}{}
\makeatletter
\@ifundefined{KOMAClassName}{% if non-KOMA class
  \IfFileExists{parskip.sty}{%
    \usepackage{parskip}
  }{% else
    \setlength{\parindent}{0pt}
    \setlength{\parskip}{6pt plus 2pt minus 1pt}}
}{% if KOMA class
  \KOMAoptions{parskip=half}}
\makeatother
\usepackage{xcolor}
\usepackage[margin=1in]{geometry}
\usepackage{color}
\usepackage{fancyvrb}
\newcommand{\VerbBar}{|}
\newcommand{\VERB}{\Verb[commandchars=\\\{\}]}
\DefineVerbatimEnvironment{Highlighting}{Verbatim}{commandchars=\\\{\}}
% Add ',fontsize=\small' for more characters per line
\usepackage{framed}
\definecolor{shadecolor}{RGB}{248,248,248}
\newenvironment{Shaded}{\begin{snugshade}}{\end{snugshade}}
\newcommand{\AlertTok}[1]{\textcolor[rgb]{0.94,0.16,0.16}{#1}}
\newcommand{\AnnotationTok}[1]{\textcolor[rgb]{0.56,0.35,0.01}{\textbf{\textit{#1}}}}
\newcommand{\AttributeTok}[1]{\textcolor[rgb]{0.13,0.29,0.53}{#1}}
\newcommand{\BaseNTok}[1]{\textcolor[rgb]{0.00,0.00,0.81}{#1}}
\newcommand{\BuiltInTok}[1]{#1}
\newcommand{\CharTok}[1]{\textcolor[rgb]{0.31,0.60,0.02}{#1}}
\newcommand{\CommentTok}[1]{\textcolor[rgb]{0.56,0.35,0.01}{\textit{#1}}}
\newcommand{\CommentVarTok}[1]{\textcolor[rgb]{0.56,0.35,0.01}{\textbf{\textit{#1}}}}
\newcommand{\ConstantTok}[1]{\textcolor[rgb]{0.56,0.35,0.01}{#1}}
\newcommand{\ControlFlowTok}[1]{\textcolor[rgb]{0.13,0.29,0.53}{\textbf{#1}}}
\newcommand{\DataTypeTok}[1]{\textcolor[rgb]{0.13,0.29,0.53}{#1}}
\newcommand{\DecValTok}[1]{\textcolor[rgb]{0.00,0.00,0.81}{#1}}
\newcommand{\DocumentationTok}[1]{\textcolor[rgb]{0.56,0.35,0.01}{\textbf{\textit{#1}}}}
\newcommand{\ErrorTok}[1]{\textcolor[rgb]{0.64,0.00,0.00}{\textbf{#1}}}
\newcommand{\ExtensionTok}[1]{#1}
\newcommand{\FloatTok}[1]{\textcolor[rgb]{0.00,0.00,0.81}{#1}}
\newcommand{\FunctionTok}[1]{\textcolor[rgb]{0.13,0.29,0.53}{\textbf{#1}}}
\newcommand{\ImportTok}[1]{#1}
\newcommand{\InformationTok}[1]{\textcolor[rgb]{0.56,0.35,0.01}{\textbf{\textit{#1}}}}
\newcommand{\KeywordTok}[1]{\textcolor[rgb]{0.13,0.29,0.53}{\textbf{#1}}}
\newcommand{\NormalTok}[1]{#1}
\newcommand{\OperatorTok}[1]{\textcolor[rgb]{0.81,0.36,0.00}{\textbf{#1}}}
\newcommand{\OtherTok}[1]{\textcolor[rgb]{0.56,0.35,0.01}{#1}}
\newcommand{\PreprocessorTok}[1]{\textcolor[rgb]{0.56,0.35,0.01}{\textit{#1}}}
\newcommand{\RegionMarkerTok}[1]{#1}
\newcommand{\SpecialCharTok}[1]{\textcolor[rgb]{0.81,0.36,0.00}{\textbf{#1}}}
\newcommand{\SpecialStringTok}[1]{\textcolor[rgb]{0.31,0.60,0.02}{#1}}
\newcommand{\StringTok}[1]{\textcolor[rgb]{0.31,0.60,0.02}{#1}}
\newcommand{\VariableTok}[1]{\textcolor[rgb]{0.00,0.00,0.00}{#1}}
\newcommand{\VerbatimStringTok}[1]{\textcolor[rgb]{0.31,0.60,0.02}{#1}}
\newcommand{\WarningTok}[1]{\textcolor[rgb]{0.56,0.35,0.01}{\textbf{\textit{#1}}}}
\usepackage{longtable,booktabs,array}
\usepackage{calc} % for calculating minipage widths
% Correct order of tables after \paragraph or \subparagraph
\usepackage{etoolbox}
\makeatletter
\patchcmd\longtable{\par}{\if@noskipsec\mbox{}\fi\par}{}{}
\makeatother
% Allow footnotes in longtable head/foot
\IfFileExists{footnotehyper.sty}{\usepackage{footnotehyper}}{\usepackage{footnote}}
\makesavenoteenv{longtable}
\usepackage{graphicx}
\makeatletter
\def\maxwidth{\ifdim\Gin@nat@width>\linewidth\linewidth\else\Gin@nat@width\fi}
\def\maxheight{\ifdim\Gin@nat@height>\textheight\textheight\else\Gin@nat@height\fi}
\makeatother
% Scale images if necessary, so that they will not overflow the page
% margins by default, and it is still possible to overwrite the defaults
% using explicit options in \includegraphics[width, height, ...]{}
\setkeys{Gin}{width=\maxwidth,height=\maxheight,keepaspectratio}
% Set default figure placement to htbp
\makeatletter
\def\fps@figure{htbp}
\makeatother
\setlength{\emergencystretch}{3em} % prevent overfull lines
\providecommand{\tightlist}{%
  \setlength{\itemsep}{0pt}\setlength{\parskip}{0pt}}
\setcounter{secnumdepth}{-\maxdimen} % remove section numbering
\usepackage{booktabs}
\usepackage{longtable}
\usepackage{array}
\usepackage{multirow}
\usepackage{wrapfig}
\usepackage{float}
\usepackage{colortbl}
\usepackage{pdflscape}
\usepackage{tabu}
\usepackage{threeparttable}
\usepackage{threeparttablex}
\usepackage[normalem]{ulem}
\usepackage{makecell}
\usepackage{xcolor}
\ifLuaTeX
  \usepackage{selnolig}  % disable illegal ligatures
\fi
\IfFileExists{bookmark.sty}{\usepackage{bookmark}}{\usepackage{hyperref}}
\IfFileExists{xurl.sty}{\usepackage{xurl}}{} % add URL line breaks if available
\urlstyle{same}
\hypersetup{
  pdftitle={Proyecto\_Entrega2\_Presentación\_Resultados},
  pdfauthor={- Luis Pedro Lira (23669) - Fernando Rocha (23501) - Juan Francisco Martínez (23617) - Luis Gilberto González (23353) - Joel Antonio Jaquez (23369)},
  hidelinks,
  pdfcreator={LaTeX via pandoc}}

\title{Proyecto\_Entrega2\_Presentación\_Resultados}
\author{- Luis Pedro Lira (23669) - Fernando Rocha (23501) - Juan
Francisco Martínez (23617) - Luis Gilberto González (23353) - Joel
Antonio Jaquez (23369)}
\date{2026-05-01}

\begin{document}
\maketitle

\begin{center}\rule{0.5\linewidth}{0.5pt}\end{center}

\hypertarget{actividad-5-conjuntos-de-entrenamiento-y-prueba}{%
\section{Actividad 5: Conjuntos de entrenamiento y
prueba}\label{actividad-5-conjuntos-de-entrenamiento-y-prueba}}

\hypertarget{muxe9todo-de-validaciuxf3n}{%
\subsection{Método de validación}\label{muxe9todo-de-validaciuxf3n}}

Dado que el conjunto de datos está agregado a nivel departamental y
cuenta con únicamente 22 observaciones, no es posible aplicar una
partición fija tradicional de entrenamiento y prueba. Con una división
estándar 80/20 el conjunto de prueba quedaría con apenas 4
departamentos, lo que haría que los resultados dependieran fuertemente
de cuáles cayeron en cada partición.

Por esta razón se utilizó \textbf{Leave-One-Out Cross Validation
(LOOCV)}. En cada iteración el modelo se entrena con 21 departamentos y
se evalúa en el departamento restante, repitiendo el proceso 22 veces
hasta que todos los departamentos hayan sido usados exactamente una vez
como prueba. Las métricas finales se calculan promediando los errores de
las 22 iteraciones, lo que produce una estimación más honesta del error
de generalización sin desperdiciar observaciones.

Como la variable respuesta es cuantitativa, \textbf{no aplica el
concepto de balance de clases}. En su lugar se analiza la presencia de
valores atípicos con alta influencia sobre el modelo.

\begin{Shaded}
\begin{Highlighting}[]
\NormalTok{n\_obs      }\OtherTok{\textless{}{-}} \FunctionTok{nrow}\NormalTok{(datos\_modelo)}
\NormalTok{porc\_train }\OtherTok{\textless{}{-}}\NormalTok{ (n\_obs }\SpecialCharTok{{-}} \DecValTok{1}\NormalTok{) }\SpecialCharTok{/}\NormalTok{ n\_obs }\SpecialCharTok{*} \DecValTok{100}
\NormalTok{porc\_test  }\OtherTok{\textless{}{-}} \DecValTok{1} \SpecialCharTok{/}\NormalTok{ n\_obs }\SpecialCharTok{*} \DecValTok{100}

\FunctionTok{tibble}\NormalTok{(}
  \StringTok{\textasciigrave{}}\AttributeTok{N observaciones}\StringTok{\textasciigrave{}}               \OtherTok{=}\NormalTok{ n\_obs,}
  \StringTok{\textasciigrave{}}\AttributeTok{\% entrenamiento por iteración}\StringTok{\textasciigrave{}} \OtherTok{=} \FunctionTok{round}\NormalTok{(porc\_train, }\DecValTok{2}\NormalTok{),}
  \StringTok{\textasciigrave{}}\AttributeTok{\% prueba por iteración}\StringTok{\textasciigrave{}}        \OtherTok{=} \FunctionTok{round}\NormalTok{(porc\_test, }\DecValTok{2}\NormalTok{)}
\NormalTok{) }\SpecialCharTok{\%\textgreater{}\%}
  \FunctionTok{kable}\NormalTok{(}\AttributeTok{caption =} \StringTok{"Proporciones de entrenamiento y prueba bajo LOOCV"}\NormalTok{)}
\end{Highlighting}
\end{Shaded}

\begin{longtable}[]{@{}
  >{\raggedleft\arraybackslash}p{(\columnwidth - 4\tabcolsep) * \real{0.2319}}
  >{\raggedleft\arraybackslash}p{(\columnwidth - 4\tabcolsep) * \real{0.4348}}
  >{\raggedleft\arraybackslash}p{(\columnwidth - 4\tabcolsep) * \real{0.3333}}@{}}
\caption{Proporciones de entrenamiento y prueba bajo
LOOCV}\tabularnewline
\toprule\noalign{}
\begin{minipage}[b]{\linewidth}\raggedleft
N observaciones
\end{minipage} & \begin{minipage}[b]{\linewidth}\raggedleft
\% entrenamiento por iteración
\end{minipage} & \begin{minipage}[b]{\linewidth}\raggedleft
\% prueba por iteración
\end{minipage} \\
\midrule\noalign{}
\endfirsthead
\toprule\noalign{}
\begin{minipage}[b]{\linewidth}\raggedleft
N observaciones
\end{minipage} & \begin{minipage}[b]{\linewidth}\raggedleft
\% entrenamiento por iteración
\end{minipage} & \begin{minipage}[b]{\linewidth}\raggedleft
\% prueba por iteración
\end{minipage} \\
\midrule\noalign{}
\endhead
\bottomrule\noalign{}
\endlastfoot
22 & 95.45 & 4.55 \\
\end{longtable}

En cada iteración se usa el \textbf{95.45\%} de los datos para
entrenamiento y el \textbf{4.55\%} para prueba.

\hypertarget{anuxe1lisis-de-atuxedpicos}{%
\subsection{Análisis de atípicos}\label{anuxe1lisis-de-atuxedpicos}}

Al tratarse de un problema de regresión con variable respuesta
cuantitativa, se analiza la distribución de \texttt{casos\_divorcio}
para identificar observaciones con comportamiento extremo que puedan
influir desproporcionadamente en los modelos.

\begin{Shaded}
\begin{Highlighting}[]
\FunctionTok{ggplot}\NormalTok{(datos\_modelo, }\FunctionTok{aes}\NormalTok{(}\AttributeTok{x =} \StringTok{""}\NormalTok{, }\AttributeTok{y =}\NormalTok{ casos\_divorcio)) }\SpecialCharTok{+}
  \FunctionTok{geom\_boxplot}\NormalTok{(}\AttributeTok{fill =} \StringTok{"lightblue"}\NormalTok{, }\AttributeTok{outlier.color =} \StringTok{"red"}\NormalTok{,}
               \AttributeTok{outlier.size =} \DecValTok{3}\NormalTok{) }\SpecialCharTok{+}
  \FunctionTok{geom\_text}\NormalTok{(}
    \AttributeTok{data =}\NormalTok{ datos\_modelo }\SpecialCharTok{\%\textgreater{}\%} \FunctionTok{filter}\NormalTok{(casos\_divorcio }\SpecialCharTok{\textgreater{}} \DecValTok{1000}\NormalTok{),}
    \FunctionTok{aes}\NormalTok{(}\AttributeTok{label =}\NormalTok{ cod\_depto),}
    \AttributeTok{hjust =} \SpecialCharTok{{-}}\FloatTok{0.5}\NormalTok{, }\AttributeTok{size =} \DecValTok{3}
\NormalTok{  ) }\SpecialCharTok{+}
  \FunctionTok{labs}\NormalTok{(}
    \AttributeTok{title =} \StringTok{"Distribución de divorcios por departamento (2022)"}\NormalTok{,}
    \AttributeTok{x     =} \StringTok{""}\NormalTok{,}
    \AttributeTok{y     =} \StringTok{"Número de divorcios"}
\NormalTok{  ) }\SpecialCharTok{+}
  \FunctionTok{theme\_minimal}\NormalTok{()}
\end{Highlighting}
\end{Shaded}

\begin{figure}
\centering
\includegraphics{Proyecto_Entrega2_Presentacion_Resultados_files/figure-latex/atipicos-boxplot-1.pdf}
\caption{Distribución de divorcios por departamento. El punto rojo
corresponde al Departamento de Guatemala, cuyo volumen es extremadamente
superior al resto del país.}
\end{figure}

\begin{Shaded}
\begin{Highlighting}[]
\NormalTok{datos\_modelo }\SpecialCharTok{\%\textgreater{}\%}
  \FunctionTok{arrange}\NormalTok{(}\FunctionTok{desc}\NormalTok{(casos\_divorcio)) }\SpecialCharTok{\%\textgreater{}\%}
  \FunctionTok{select}\NormalTok{(cod\_depto, casos\_divorcio) }\SpecialCharTok{\%\textgreater{}\%}
  \FunctionTok{mutate}\NormalTok{(}
    \AttributeTok{z\_score =} \FunctionTok{round}\NormalTok{(}
\NormalTok{      (casos\_divorcio }\SpecialCharTok{{-}} \FunctionTok{mean}\NormalTok{(casos\_divorcio)) }\SpecialCharTok{/} \FunctionTok{sd}\NormalTok{(casos\_divorcio), }\DecValTok{2}
\NormalTok{    )}
\NormalTok{  ) }\SpecialCharTok{\%\textgreater{}\%}
  \FunctionTok{kable}\NormalTok{(}\AttributeTok{caption =} \StringTok{"Departamentos ordenados por divorcios con z{-}score"}\NormalTok{)}
\end{Highlighting}
\end{Shaded}

\begin{longtable}[]{@{}rrr@{}}
\caption{Departamentos ordenados por divorcios con
z-score}\tabularnewline
\toprule\noalign{}
cod\_depto & casos\_divorcio & z\_score \\
\midrule\noalign{}
\endfirsthead
\toprule\noalign{}
cod\_depto & casos\_divorcio & z\_score \\
\midrule\noalign{}
\endhead
\bottomrule\noalign{}
\endlastfoot
1 & 3295 & 4.34 \\
9 & 794 & 0.52 \\
12 & 565 & 0.17 \\
5 & 497 & 0.07 \\
10 & 485 & 0.05 \\
13 & 450 & 0.00 \\
22 & 402 & -0.08 \\
11 & 331 & -0.19 \\
17 & 326 & -0.19 \\
6 & 316 & -0.21 \\
14 & 314 & -0.21 \\
18 & 286 & -0.25 \\
21 & 238 & -0.33 \\
20 & 212 & -0.37 \\
4 & 210 & -0.37 \\
8 & 201 & -0.38 \\
16 & 198 & -0.39 \\
19 & 195 & -0.39 \\
2 & 179 & -0.42 \\
3 & 174 & -0.42 \\
15 & 145 & -0.47 \\
7 & 137 & -0.48 \\
\end{longtable}

El Departamento de Guatemala (cod\_depto = 1) presenta un z-score muy
superior al resto, confirmando su condición de \textbf{outlier
estructural extremo}. Su volumen de divorcios es entre 6 y 8 veces mayor
que el promedio del resto de departamentos, lo que refleja su
desproporción poblacional y su concentración de servicios judiciales
respecto al interior del país.

Este departamento no se elimina del análisis porque es una observación
real y relevante para el problema. En su lugar se controla mediante la
variable dummy \texttt{es\_capital}, que se detalla en la siguiente
sección.

\begin{center}\rule{0.5\linewidth}{0.5pt}\end{center}

\hypertarget{actividad-6-preprocesamiento-y-preparaciuxf3n-de-datos}{%
\section{Actividad 6: Preprocesamiento y preparación de
datos}\label{actividad-6-preprocesamiento-y-preparaciuxf3n-de-datos}}

Esta sección describe todas las transformaciones aplicadas a los datos
antes de entrenar los modelos. El resultado final es el objeto
\texttt{datos\_modelo}, que será utilizado por todos los algoritmos en
las actividades 7 y 8.

\hypertarget{carga-y-filtro-por-auxf1o}{%
\subsection{Carga y filtro por año}\label{carga-y-filtro-por-auxf1o}}

Los archivos fuente del INE contienen registros del período 2012 a 2022.
Se filtra únicamente el año \textbf{2022} para garantizar un análisis de
corte transversal limpio, donde cada departamento representa una
observación bajo condiciones temporales homogéneas. Mezclar múltiples
años sin modelar la estructura temporal produciría conteos inflados sin
interpretación válida para un modelo de regresión de corte transversal.

\begin{Shaded}
\begin{Highlighting}[]
\FunctionTok{cat}\NormalTok{(}\StringTok{"Registros de divorcios (2022):"}\NormalTok{, }\FunctionTok{nrow}\NormalTok{(divorcios\_raw), }\StringTok{"}\SpecialCharTok{\textbackslash{}n}\StringTok{"}\NormalTok{)}
\end{Highlighting}
\end{Shaded}

\begin{verbatim}
## Registros de divorcios (2022): 9950
\end{verbatim}

\begin{Shaded}
\begin{Highlighting}[]
\FunctionTok{cat}\NormalTok{(}\StringTok{"Registros de violencia (2022):"}\NormalTok{, }\FunctionTok{nrow}\NormalTok{(violencia\_raw), }\StringTok{"}\SpecialCharTok{\textbackslash{}n}\StringTok{"}\NormalTok{)}
\end{Highlighting}
\end{Shaded}

\begin{verbatim}
## Registros de violencia (2022): 37194
\end{verbatim}

\hypertarget{limpieza-y-recodificaciuxf3n}{%
\subsection{Limpieza y
recodificación}\label{limpieza-y-recodificaciuxf3n}}

Los archivos del INE utilizan los códigos 99 y 999 para representar
información no especificada. Estos valores se transforman en \texttt{NA}
para evitar que distorsionen los cálculos. Adicionalmente, los registros
con código departamental 99 se excluyen porque no pueden asignarse a
ninguna unidad geográfica válida.

\hypertarget{agregaciuxf3n-departamental-y-selecciuxf3n-de-covariables}{%
\subsection{Agregación departamental y selección de
covariables}\label{agregaciuxf3n-departamental-y-selecciuxf3n-de-covariables}}

Los microdatos individuales se agrupan por departamento. Se calculan
tres covariables a partir del dataset de violencia:

\begin{itemize}
\tightlist
\item
  \textbf{\texttt{casos\_violencia}}: conteo total de casos de violencia
  intrafamiliar.
\item
  \textbf{\texttt{edad\_victima\_prom}}: edad promedio de la víctima
  como proxy del perfil etario departamental.
\item
  \textbf{\texttt{porc\_urbano}}: proporción de casos en área urbana
  como proxy de urbanización.
\end{itemize}

La variable \texttt{hijos\_prom} fue evaluada y descartada por presentar
una correlación prácticamente nula con la variable respuesta (r =
0.095), lo que indica que no aporta capacidad predictiva al modelo.

\begin{Shaded}
\begin{Highlighting}[]
\FunctionTok{cat}\NormalTok{(}\StringTok{"Departamentos en violencia:"}\NormalTok{, }\FunctionTok{nrow}\NormalTok{(violencia\_dep), }\StringTok{"}\SpecialCharTok{\textbackslash{}n}\StringTok{"}\NormalTok{)}
\end{Highlighting}
\end{Shaded}

\begin{verbatim}
## Departamentos en violencia: 22
\end{verbatim}

\begin{Shaded}
\begin{Highlighting}[]
\FunctionTok{cat}\NormalTok{(}\StringTok{"Departamentos en divorcios:"}\NormalTok{, }\FunctionTok{nrow}\NormalTok{(divorcios\_dep), }\StringTok{"}\SpecialCharTok{\textbackslash{}n}\StringTok{"}\NormalTok{)}
\end{Highlighting}
\end{Shaded}

\begin{verbatim}
## Departamentos en divorcios: 22
\end{verbatim}

\hypertarget{transformaciones-aplicadas-al-dataset-final}{%
\subsection{Transformaciones aplicadas al dataset
final}\label{transformaciones-aplicadas-al-dataset-final}}

Una vez construido el dataset con las variables base, se aplicaron dos
transformaciones adicionales:

\textbf{Estandarización z-score.} Las variables continuas
(\texttt{casos\_violencia}, \texttt{edad\_victima\_prom},
\texttt{porc\_urbano}) se transformaron a z-score, restando la media y
dividiendo por la desviación estándar. Esto permite comparar
directamente la magnitud de los coeficientes entre predictores con
escalas distintas y mejora la estabilidad numérica de los algoritmos.

\textbf{Variable dummy \texttt{es\_capital}.} El análisis exploratorio
identificó al Departamento de Guatemala como un outlier estructural
extremo. Esta variable binaria vale 1 para ese departamento y 0 para los
demás, permitiendo que los modelos capturen su diferencia estructural
sin eliminarlo del análisis.

\hypertarget{dataset-final-de-modelado}{%
\subsection{Dataset final de modelado}\label{dataset-final-de-modelado}}

\begin{Shaded}
\begin{Highlighting}[]
\NormalTok{datos\_modelo }\SpecialCharTok{\%\textgreater{}\%}
  \FunctionTok{select}\NormalTok{(cod\_depto, casos\_divorcio, casos\_violencia,}
\NormalTok{         edad\_victima\_prom, porc\_urbano, es\_capital) }\SpecialCharTok{\%\textgreater{}\%}
  \FunctionTok{kable}\NormalTok{(}\AttributeTok{digits =} \DecValTok{2}\NormalTok{,}
        \AttributeTok{caption =} \StringTok{"Dataset de modelado: variables originales por departamento (2022)"}\NormalTok{)}
\end{Highlighting}
\end{Shaded}

\begin{longtable}[]{@{}
  >{\raggedleft\arraybackslash}p{(\columnwidth - 10\tabcolsep) * \real{0.1220}}
  >{\raggedleft\arraybackslash}p{(\columnwidth - 10\tabcolsep) * \real{0.1829}}
  >{\raggedleft\arraybackslash}p{(\columnwidth - 10\tabcolsep) * \real{0.1951}}
  >{\raggedleft\arraybackslash}p{(\columnwidth - 10\tabcolsep) * \real{0.2195}}
  >{\raggedleft\arraybackslash}p{(\columnwidth - 10\tabcolsep) * \real{0.1463}}
  >{\raggedleft\arraybackslash}p{(\columnwidth - 10\tabcolsep) * \real{0.1341}}@{}}
\caption{Dataset de modelado: variables originales por departamento
(2022)}\tabularnewline
\toprule\noalign{}
\begin{minipage}[b]{\linewidth}\raggedleft
cod\_depto
\end{minipage} & \begin{minipage}[b]{\linewidth}\raggedleft
casos\_divorcio
\end{minipage} & \begin{minipage}[b]{\linewidth}\raggedleft
casos\_violencia
\end{minipage} & \begin{minipage}[b]{\linewidth}\raggedleft
edad\_victima\_prom
\end{minipage} & \begin{minipage}[b]{\linewidth}\raggedleft
porc\_urbano
\end{minipage} & \begin{minipage}[b]{\linewidth}\raggedleft
es\_capital
\end{minipage} \\
\midrule\noalign{}
\endfirsthead
\toprule\noalign{}
\begin{minipage}[b]{\linewidth}\raggedleft
cod\_depto
\end{minipage} & \begin{minipage}[b]{\linewidth}\raggedleft
casos\_divorcio
\end{minipage} & \begin{minipage}[b]{\linewidth}\raggedleft
casos\_violencia
\end{minipage} & \begin{minipage}[b]{\linewidth}\raggedleft
edad\_victima\_prom
\end{minipage} & \begin{minipage}[b]{\linewidth}\raggedleft
porc\_urbano
\end{minipage} & \begin{minipage}[b]{\linewidth}\raggedleft
es\_capital
\end{minipage} \\
\midrule\noalign{}
\endhead
\bottomrule\noalign{}
\endlastfoot
1 & 3295 & 8268 & 39.46 & 0.88 & 1 \\
2 & 179 & 1122 & 38.78 & 0.40 & 0 \\
3 & 174 & 1689 & 36.18 & 0.73 & 0 \\
4 & 210 & 2169 & 34.47 & 0.66 & 0 \\
5 & 497 & 1745 & 33.18 & 0.70 & 0 \\
6 & 316 & 1126 & 35.23 & 0.32 & 0 \\
7 & 137 & 355 & 40.12 & 0.32 & 0 \\
8 & 201 & 462 & 35.41 & 0.26 & 0 \\
9 & 794 & 2106 & 36.77 & 0.61 & 0 \\
10 & 485 & 1765 & 35.63 & 0.45 & 0 \\
11 & 331 & 1575 & 32.08 & 0.37 & 0 \\
12 & 565 & 1703 & 37.37 & 0.29 & 0 \\
13 & 450 & 999 & 36.04 & 0.34 & 0 \\
14 & 314 & 629 & 36.01 & 0.35 & 0 \\
15 & 145 & 1331 & 37.43 & 0.37 & 0 \\
16 & 198 & 4561 & 33.86 & 0.43 & 0 \\
17 & 326 & 1043 & 34.37 & 0.80 & 0 \\
18 & 286 & 918 & 35.26 & 0.51 & 0 \\
19 & 195 & 903 & 38.00 & 0.49 & 0 \\
20 & 212 & 475 & 37.07 & 0.56 & 0 \\
21 & 238 & 456 & 35.27 & 0.47 & 0 \\
22 & 402 & 1245 & 38.21 & 0.26 & 0 \\
\end{longtable}

\begin{Shaded}
\begin{Highlighting}[]
\NormalTok{datos\_modelo }\SpecialCharTok{\%\textgreater{}\%}
  \FunctionTok{select}\NormalTok{(casos\_divorcio, casos\_violencia,}
\NormalTok{         edad\_victima\_prom, porc\_urbano) }\SpecialCharTok{\%\textgreater{}\%}
  \FunctionTok{summary}\NormalTok{()}
\end{Highlighting}
\end{Shaded}

\begin{verbatim}
##  casos_divorcio   casos_violencia  edad_victima_prom  porc_urbano    
##  Min.   : 137.0   Min.   : 355.0   Min.   :32.08     Min.   :0.2602  
##  1st Qu.: 198.8   1st Qu.: 906.8   1st Qu.:35.24     1st Qu.:0.3457  
##  Median : 300.0   Median :1185.5   Median :36.03     Median :0.4421  
##  Mean   : 452.3   Mean   :1665.7   Mean   :36.19     Mean   :0.4805  
##  3rd Qu.: 438.0   3rd Qu.:1734.5   3rd Qu.:37.41     3rd Qu.:0.5969  
##  Max.   :3295.0   Max.   :8268.0   Max.   :40.12     Max.   :0.8805
\end{verbatim}

El dataset final cuenta con \textbf{22 observaciones} una por
departamento y las columnas estandarizadas (\texttt{\_z}) listas para
ser usadas directamente por todos los modelos. No hay valores faltantes
en ninguna de las covariables, por lo que no fue necesario aplicar
imputación.

\hypertarget{funciuxf3n-de-muxe9tricas}{%
\subsection{Función de métricas}\label{funciuxf3n-de-muxe9tricas}}

Para mantener consistencia en la evaluación de todos los algoritmos, se
define una única función de métricas que será utilizada por los cuatro
modelos en las actividades siguientes:

\begin{itemize}
\tightlist
\item
  \textbf{MAE}: error absoluto promedio en número de divorcios.
\item
  \textbf{RMSE}: penaliza más los errores grandes, útil para detectar
  fallos en departamentos extremos.
\item
  \textbf{MAPE}: error porcentual promedio, facilita la interpretación
  relativa.
\end{itemize}

Para los modelos estadísticos (Poisson y Binomial Negativa) se añaden
adicionalmente \textbf{AIC} y \textbf{BIC} como criterios de selección
entre variantes.

\begin{center}\rule{0.5\linewidth}{0.5pt}\end{center}

\hypertarget{actividad-7-creaciuxf3n-de-modelos-por-algoritmos}{%
\section{Actividad 7: Creación de modelos por
algoritmos}\label{actividad-7-creaciuxf3n-de-modelos-por-algoritmos}}

\hypertarget{regresiuxf3n-binomial-negativa}{%
\subsection{7.1 Regresión Binomial
Negativa}\label{regresiuxf3n-binomial-negativa}}

La Regresión Binomial Negativa es el modelo principal del análisis. Se
implementan cuatro variantes con distintas combinaciones de predictores
para identificar cuál logra el mejor equilibrio entre ajuste y
parsimonia. Dado que el conjunto de datos tiene únicamente 22
observaciones, se evalúa cada variante mediante LOOCV en lugar de un
split tradicional.

\begin{Shaded}
\begin{Highlighting}[]
\NormalTok{loo\_nb }\OtherTok{\textless{}{-}} \ControlFlowTok{function}\NormalTok{(formula\_modelo, datos) \{}
\NormalTok{  n            }\OtherTok{\textless{}{-}} \FunctionTok{nrow}\NormalTok{(datos)}
\NormalTok{  predicciones }\OtherTok{\textless{}{-}} \FunctionTok{numeric}\NormalTok{(n)}

  \ControlFlowTok{for}\NormalTok{ (i }\ControlFlowTok{in} \FunctionTok{seq\_len}\NormalTok{(n)) \{}
\NormalTok{    entrenamiento   }\OtherTok{\textless{}{-}}\NormalTok{ datos[}\SpecialCharTok{{-}}\NormalTok{i, , drop }\OtherTok{=} \ConstantTok{FALSE}\NormalTok{]}
\NormalTok{    prueba          }\OtherTok{\textless{}{-}}\NormalTok{ datos[i, , drop }\OtherTok{=} \ConstantTok{FALSE}\NormalTok{]}
\NormalTok{    modelo          }\OtherTok{\textless{}{-}}\NormalTok{ MASS}\SpecialCharTok{::}\FunctionTok{glm.nb}\NormalTok{(formula\_modelo, }\AttributeTok{data =}\NormalTok{ entrenamiento)}
\NormalTok{    predicciones[i] }\OtherTok{\textless{}{-}} \FunctionTok{predict}\NormalTok{(modelo, }\AttributeTok{newdata =}\NormalTok{ prueba, }\AttributeTok{type =} \StringTok{"response"}\NormalTok{)}
\NormalTok{  \}}

  \FunctionTok{list}\NormalTok{(}
    \AttributeTok{metricas     =} \FunctionTok{calcular\_metricas}\NormalTok{(datos}\SpecialCharTok{$}\NormalTok{casos\_divorcio, predicciones),}
    \AttributeTok{predicciones =}\NormalTok{ predicciones}
\NormalTok{  )}
\NormalTok{\}}
\end{Highlighting}
\end{Shaded}

Se define la función loo\_nb que implementa LOOCV para cualquier fórmula
de Binomial Negativa. En cada iteración entrena con 21 departamentos y
predice el restante, devolviendo las métricas MAE, RMSE y MAPE junto con
el vector de predicciones.

\begin{Shaded}
\begin{Highlighting}[]
\NormalTok{formula\_nb\_1 }\OtherTok{\textless{}{-}}\NormalTok{ casos\_divorcio }\SpecialCharTok{\textasciitilde{}}\NormalTok{ casos\_violencia\_z }\SpecialCharTok{+}\NormalTok{ es\_capital}

\NormalTok{formula\_nb\_2 }\OtherTok{\textless{}{-}}\NormalTok{ casos\_divorcio }\SpecialCharTok{\textasciitilde{}}\NormalTok{ casos\_violencia\_z }\SpecialCharTok{+}\NormalTok{ es\_capital }\SpecialCharTok{+}
\NormalTok{                porc\_urbano\_z}

\NormalTok{formula\_nb\_3 }\OtherTok{\textless{}{-}}\NormalTok{ casos\_divorcio }\SpecialCharTok{\textasciitilde{}}\NormalTok{ casos\_violencia\_z }\SpecialCharTok{*}\NormalTok{ es\_capital}

\NormalTok{formula\_nb\_4 }\OtherTok{\textless{}{-}}\NormalTok{ casos\_divorcio }\SpecialCharTok{\textasciitilde{}}\NormalTok{ casos\_violencia\_z }\SpecialCharTok{+}\NormalTok{ es\_capital }\SpecialCharTok{+}
\NormalTok{                porc\_urbano\_z }\SpecialCharTok{+}\NormalTok{ edad\_victima\_prom\_z}
\end{Highlighting}
\end{Shaded}

Se definen cuatro fórmulas que varían en complejidad. NB\_1 es el modelo
base con los dos predictores más importantes identificados en el
análisis exploratorio. NB\_2 agrega la proporción de urbanización. NB\_3
introduce un término de interacción entre violencia y la dummy de
capital para capturar si el efecto de la violencia sobre los divorcios
es distinto en Guatemala respecto al interior. NB\_4 es el modelo más
complejo, incluyendo además la edad promedio de la víctima.

\begin{Shaded}
\begin{Highlighting}[]
\NormalTok{modelo\_nb\_1 }\OtherTok{\textless{}{-}}\NormalTok{ MASS}\SpecialCharTok{::}\FunctionTok{glm.nb}\NormalTok{(formula\_nb\_1, }\AttributeTok{data =}\NormalTok{ datos\_modelo)}
\NormalTok{modelo\_nb\_2 }\OtherTok{\textless{}{-}}\NormalTok{ MASS}\SpecialCharTok{::}\FunctionTok{glm.nb}\NormalTok{(formula\_nb\_2, }\AttributeTok{data =}\NormalTok{ datos\_modelo)}
\NormalTok{modelo\_nb\_3 }\OtherTok{\textless{}{-}}\NormalTok{ MASS}\SpecialCharTok{::}\FunctionTok{glm.nb}\NormalTok{(formula\_nb\_3, }\AttributeTok{data =}\NormalTok{ datos\_modelo)}
\NormalTok{modelo\_nb\_4 }\OtherTok{\textless{}{-}}\NormalTok{ MASS}\SpecialCharTok{::}\FunctionTok{glm.nb}\NormalTok{(formula\_nb\_4, }\AttributeTok{data =}\NormalTok{ datos\_modelo)}
\end{Highlighting}
\end{Shaded}

Se ajustan los cuatro modelos sobre el dataset completo usando
MASS::glm.nb. Estos objetos se usan para extraer AIC y BIC, que
complementan las métricas LOOCV en la comparación.

\begin{Shaded}
\begin{Highlighting}[]
\NormalTok{loo\_nb\_1 }\OtherTok{\textless{}{-}} \FunctionTok{loo\_nb}\NormalTok{(formula\_nb\_1, datos\_modelo)}
\NormalTok{loo\_nb\_2 }\OtherTok{\textless{}{-}} \FunctionTok{loo\_nb}\NormalTok{(formula\_nb\_2, datos\_modelo)}
\NormalTok{loo\_nb\_3 }\OtherTok{\textless{}{-}} \FunctionTok{loo\_nb}\NormalTok{(formula\_nb\_3, datos\_modelo)}
\NormalTok{loo\_nb\_4 }\OtherTok{\textless{}{-}} \FunctionTok{loo\_nb}\NormalTok{(formula\_nb\_4, datos\_modelo)}

\NormalTok{pred\_loo\_nb }\OtherTok{\textless{}{-}} \FunctionTok{list}\NormalTok{(}
  \AttributeTok{NB\_1 =}\NormalTok{ loo\_nb\_1}\SpecialCharTok{$}\NormalTok{predicciones,}
  \AttributeTok{NB\_2 =}\NormalTok{ loo\_nb\_2}\SpecialCharTok{$}\NormalTok{predicciones,}
  \AttributeTok{NB\_3 =}\NormalTok{ loo\_nb\_3}\SpecialCharTok{$}\NormalTok{predicciones,}
  \AttributeTok{NB\_4 =}\NormalTok{ loo\_nb\_4}\SpecialCharTok{$}\NormalTok{predicciones}
\NormalTok{)}
\end{Highlighting}
\end{Shaded}

Se corre LOOCV para cada variante. Las predicciones se guardan en la
lista pred\_loo\_nb para usarlas en la comparación final del documento.

\begin{Shaded}
\begin{Highlighting}[]
\NormalTok{metricas\_nb }\OtherTok{\textless{}{-}} \FunctionTok{bind\_rows}\NormalTok{(}
\NormalTok{  loo\_nb\_1}\SpecialCharTok{$}\NormalTok{metricas }\SpecialCharTok{\%\textgreater{}\%} \FunctionTok{mutate}\NormalTok{(}\AttributeTok{modelo =} \StringTok{"NB\_1"}\NormalTok{,}
    \AttributeTok{formula =} \FunctionTok{paste}\NormalTok{(}\FunctionTok{deparse}\NormalTok{(formula\_nb\_1), }\AttributeTok{collapse =} \StringTok{" "}\NormalTok{)),}
\NormalTok{  loo\_nb\_2}\SpecialCharTok{$}\NormalTok{metricas }\SpecialCharTok{\%\textgreater{}\%} \FunctionTok{mutate}\NormalTok{(}\AttributeTok{modelo =} \StringTok{"NB\_2"}\NormalTok{,}
    \AttributeTok{formula =} \FunctionTok{paste}\NormalTok{(}\FunctionTok{deparse}\NormalTok{(formula\_nb\_2), }\AttributeTok{collapse =} \StringTok{" "}\NormalTok{)),}
\NormalTok{  loo\_nb\_3}\SpecialCharTok{$}\NormalTok{metricas }\SpecialCharTok{\%\textgreater{}\%} \FunctionTok{mutate}\NormalTok{(}\AttributeTok{modelo =} \StringTok{"NB\_3"}\NormalTok{,}
    \AttributeTok{formula =} \FunctionTok{paste}\NormalTok{(}\FunctionTok{deparse}\NormalTok{(formula\_nb\_3), }\AttributeTok{collapse =} \StringTok{" "}\NormalTok{)),}
\NormalTok{  loo\_nb\_4}\SpecialCharTok{$}\NormalTok{metricas }\SpecialCharTok{\%\textgreater{}\%} \FunctionTok{mutate}\NormalTok{(}\AttributeTok{modelo =} \StringTok{"NB\_4"}\NormalTok{,}
    \AttributeTok{formula =} \FunctionTok{paste}\NormalTok{(}\FunctionTok{deparse}\NormalTok{(formula\_nb\_4), }\AttributeTok{collapse =} \StringTok{" "}\NormalTok{))}
\NormalTok{) }\SpecialCharTok{\%\textgreater{}\%}
  \FunctionTok{mutate}\NormalTok{(}
    \AttributeTok{AIC       =} \FunctionTok{c}\NormalTok{(}\FunctionTok{AIC}\NormalTok{(modelo\_nb\_1), }\FunctionTok{AIC}\NormalTok{(modelo\_nb\_2),}
                  \FunctionTok{AIC}\NormalTok{(modelo\_nb\_3), }\FunctionTok{AIC}\NormalTok{(modelo\_nb\_4)),}
    \AttributeTok{BIC       =} \FunctionTok{c}\NormalTok{(}\FunctionTok{BIC}\NormalTok{(modelo\_nb\_1), }\FunctionTok{BIC}\NormalTok{(modelo\_nb\_2),}
                  \FunctionTok{BIC}\NormalTok{(modelo\_nb\_3), }\FunctionTok{BIC}\NormalTok{(modelo\_nb\_4)),}
    \AttributeTok{algoritmo =} \StringTok{"Regresión Binomial Negativa"}\NormalTok{,}
    \AttributeTok{rol       =} \StringTok{"Modelo principal"}
\NormalTok{  ) }\SpecialCharTok{\%\textgreater{}\%}
  \FunctionTok{select}\NormalTok{(algoritmo, rol, modelo, formula, AIC, BIC, MAE, RMSE, MAPE)}

\NormalTok{metricas\_nb }\SpecialCharTok{\%\textgreater{}\%}
  \FunctionTok{arrange}\NormalTok{(MAE, RMSE, AIC) }\SpecialCharTok{\%\textgreater{}\%}
  \FunctionTok{kable}\NormalTok{(}\AttributeTok{digits =} \DecValTok{3}\NormalTok{,}
        \AttributeTok{caption =} \StringTok{"Comparación de variantes — Regresión Binomial Negativa"}\NormalTok{)}
\end{Highlighting}
\end{Shaded}

\begin{longtable}[]{@{}
  >{\raggedright\arraybackslash}p{(\columnwidth - 16\tabcolsep) * \real{0.1538}}
  >{\raggedright\arraybackslash}p{(\columnwidth - 16\tabcolsep) * \real{0.0934}}
  >{\raggedright\arraybackslash}p{(\columnwidth - 16\tabcolsep) * \real{0.0385}}
  >{\raggedright\arraybackslash}p{(\columnwidth - 16\tabcolsep) * \real{0.5000}}
  >{\raggedleft\arraybackslash}p{(\columnwidth - 16\tabcolsep) * \real{0.0440}}
  >{\raggedleft\arraybackslash}p{(\columnwidth - 16\tabcolsep) * \real{0.0440}}
  >{\raggedleft\arraybackslash}p{(\columnwidth - 16\tabcolsep) * \real{0.0440}}
  >{\raggedleft\arraybackslash}p{(\columnwidth - 16\tabcolsep) * \real{0.0440}}
  >{\raggedleft\arraybackslash}p{(\columnwidth - 16\tabcolsep) * \real{0.0385}}@{}}
\caption{Comparación de variantes --- Regresión Binomial
Negativa}\tabularnewline
\toprule\noalign{}
\begin{minipage}[b]{\linewidth}\raggedright
algoritmo
\end{minipage} & \begin{minipage}[b]{\linewidth}\raggedright
rol
\end{minipage} & \begin{minipage}[b]{\linewidth}\raggedright
modelo
\end{minipage} & \begin{minipage}[b]{\linewidth}\raggedright
formula
\end{minipage} & \begin{minipage}[b]{\linewidth}\raggedleft
AIC
\end{minipage} & \begin{minipage}[b]{\linewidth}\raggedleft
BIC
\end{minipage} & \begin{minipage}[b]{\linewidth}\raggedleft
MAE
\end{minipage} & \begin{minipage}[b]{\linewidth}\raggedleft
RMSE
\end{minipage} & \begin{minipage}[b]{\linewidth}\raggedleft
MAPE
\end{minipage} \\
\midrule\noalign{}
\endfirsthead
\toprule\noalign{}
\begin{minipage}[b]{\linewidth}\raggedright
algoritmo
\end{minipage} & \begin{minipage}[b]{\linewidth}\raggedright
rol
\end{minipage} & \begin{minipage}[b]{\linewidth}\raggedright
modelo
\end{minipage} & \begin{minipage}[b]{\linewidth}\raggedright
formula
\end{minipage} & \begin{minipage}[b]{\linewidth}\raggedleft
AIC
\end{minipage} & \begin{minipage}[b]{\linewidth}\raggedleft
BIC
\end{minipage} & \begin{minipage}[b]{\linewidth}\raggedleft
MAE
\end{minipage} & \begin{minipage}[b]{\linewidth}\raggedleft
RMSE
\end{minipage} & \begin{minipage}[b]{\linewidth}\raggedleft
MAPE
\end{minipage} \\
\midrule\noalign{}
\endhead
\bottomrule\noalign{}
\endlastfoot
Regresión Binomial Negativa & Modelo principal & NB\_1 & casos\_divorcio
\textasciitilde{} casos\_violencia\_z + es\_capital & 289.545 & 293.909
& 290.177 & 616.737 & 78.858 \\
Regresión Binomial Negativa & Modelo principal & NB\_3 & casos\_divorcio
\textasciitilde{} casos\_violencia\_z * es\_capital & 289.545 & 293.909
& 290.177 & 616.737 & 78.858 \\
Regresión Binomial Negativa & Modelo principal & NB\_2 & casos\_divorcio
\textasciitilde{} casos\_violencia\_z + es\_capital + porc\_urbano\_z &
291.530 & 296.986 & 307.394 & 647.216 & 86.598 \\
Regresión Binomial Negativa & Modelo principal & NB\_4 & casos\_divorcio
\textasciitilde{} casos\_violencia\_z + es\_capital + porc\_urbano\_z +
edad\_victima\_prom\_z & 293.256 & 299.802 & 329.139 & 703.622 &
91.208 \\
\end{longtable}

La tabla compara las cuatro variantes del modelo Binomial Negativo
ordenadas de menor a mayor MAE bajo LOOCV. Los modelos \textbf{NB\_1} y
\textbf{NB\_3} empatan en todas las métricas (AIC = 289.545, MAE =
290.177, RMSE = 616.737), lo que indica que agregar el término de
interacción entre violencia y capital en NB\_3 no aporta mejora real
sobre el modelo base NB\_1.

Los modelos NB\_2 y NB\_4, que incluyen \texttt{porc\_urbano\_z} y
\texttt{edad\_victima\_prom\_z}, obtienen peores métricas en todos los
criterios. Esto confirma lo que anticipaba la exploración previa: con n
= 22, cada predictor adicional consume grados de libertad que los datos
no tienen capacidad de sostener, resultando en peor generalización según
el LOOCV. Este es exactamente el fenómeno de \textbf{sobreajuste} que se
buscaba evitar.

\textbf{Se selecciona NB\_1 como modelo final} por ser el más
parsimonioso con el menor MAE y AIC. Su fórmula es:

\[\log(\hat{\mu}) = \beta_0 + \beta_1 \cdot \text{casos\_violencia\_z} +
\beta_2 \cdot \text{es\_capital}\]

Esta decisión prioriza la interpretabilidad y la capacidad de
generalización sobre la complejidad, criterios especialmente relevantes
cuando el tamaño muestral es reducido.

\begin{Shaded}
\begin{Highlighting}[]
\NormalTok{mejor\_nb\_nombre }\OtherTok{\textless{}{-}}\NormalTok{ metricas\_nb }\SpecialCharTok{\%\textgreater{}\%}
  \FunctionTok{arrange}\NormalTok{(MAE, RMSE, AIC) }\SpecialCharTok{\%\textgreater{}\%}
  \FunctionTok{slice}\NormalTok{(}\DecValTok{1}\NormalTok{) }\SpecialCharTok{\%\textgreater{}\%}
  \FunctionTok{pull}\NormalTok{(modelo)}

\NormalTok{modelo\_nb\_final }\OtherTok{\textless{}{-}} \ControlFlowTok{switch}\NormalTok{(}
\NormalTok{  mejor\_nb\_nombre,}
  \StringTok{"NB\_1"} \OtherTok{=}\NormalTok{ modelo\_nb\_1,}
  \StringTok{"NB\_2"} \OtherTok{=}\NormalTok{ modelo\_nb\_2,}
  \StringTok{"NB\_3"} \OtherTok{=}\NormalTok{ modelo\_nb\_3,}
  \StringTok{"NB\_4"} \OtherTok{=}\NormalTok{ modelo\_nb\_4}
\NormalTok{)}

\NormalTok{formula\_nb\_final  }\OtherTok{\textless{}{-}} \FunctionTok{formula}\NormalTok{(modelo\_nb\_final)}
\NormalTok{pred\_nb\_final     }\OtherTok{\textless{}{-}} \FunctionTok{predict}\NormalTok{(modelo\_nb\_final,}
                             \AttributeTok{newdata =}\NormalTok{ datos\_modelo,}
                             \AttributeTok{type    =} \StringTok{"response"}\NormalTok{)}
\NormalTok{metricas\_nb\_final }\OtherTok{\textless{}{-}} \FunctionTok{calcular\_metricas}\NormalTok{(datos\_modelo}\SpecialCharTok{$}\NormalTok{casos\_divorcio,}
\NormalTok{                                       pred\_nb\_final)}
\end{Highlighting}
\end{Shaded}

\begin{Shaded}
\begin{Highlighting}[]
\NormalTok{broom}\SpecialCharTok{::}\FunctionTok{tidy}\NormalTok{(modelo\_nb\_final, }\AttributeTok{conf.int =} \ConstantTok{TRUE}\NormalTok{, }\AttributeTok{exponentiate =} \ConstantTok{TRUE}\NormalTok{) }\SpecialCharTok{\%\textgreater{}\%}
  \FunctionTok{kable}\NormalTok{(}\AttributeTok{digits =} \DecValTok{4}\NormalTok{,}
        \AttributeTok{caption =} \StringTok{"Coeficientes del modelo Binomial Negativo final (IRR con IC 95\%)"}\NormalTok{)}
\end{Highlighting}
\end{Shaded}

\begin{longtable}[]{@{}
  >{\raggedright\arraybackslash}p{(\columnwidth - 12\tabcolsep) * \real{0.2432}}
  >{\raggedleft\arraybackslash}p{(\columnwidth - 12\tabcolsep) * \real{0.1216}}
  >{\raggedleft\arraybackslash}p{(\columnwidth - 12\tabcolsep) * \real{0.1351}}
  >{\raggedleft\arraybackslash}p{(\columnwidth - 12\tabcolsep) * \real{0.1351}}
  >{\raggedleft\arraybackslash}p{(\columnwidth - 12\tabcolsep) * \real{0.1081}}
  >{\raggedleft\arraybackslash}p{(\columnwidth - 12\tabcolsep) * \real{0.1216}}
  >{\raggedleft\arraybackslash}p{(\columnwidth - 12\tabcolsep) * \real{0.1351}}@{}}
\caption{Coeficientes del modelo Binomial Negativo final (IRR con IC
95\%)}\tabularnewline
\toprule\noalign{}
\begin{minipage}[b]{\linewidth}\raggedright
term
\end{minipage} & \begin{minipage}[b]{\linewidth}\raggedleft
estimate
\end{minipage} & \begin{minipage}[b]{\linewidth}\raggedleft
std.error
\end{minipage} & \begin{minipage}[b]{\linewidth}\raggedleft
statistic
\end{minipage} & \begin{minipage}[b]{\linewidth}\raggedleft
p.value
\end{minipage} & \begin{minipage}[b]{\linewidth}\raggedleft
conf.low
\end{minipage} & \begin{minipage}[b]{\linewidth}\raggedleft
conf.high
\end{minipage} \\
\midrule\noalign{}
\endfirsthead
\toprule\noalign{}
\begin{minipage}[b]{\linewidth}\raggedright
term
\end{minipage} & \begin{minipage}[b]{\linewidth}\raggedleft
estimate
\end{minipage} & \begin{minipage}[b]{\linewidth}\raggedleft
std.error
\end{minipage} & \begin{minipage}[b]{\linewidth}\raggedleft
statistic
\end{minipage} & \begin{minipage}[b]{\linewidth}\raggedleft
p.value
\end{minipage} & \begin{minipage}[b]{\linewidth}\raggedleft
conf.low
\end{minipage} & \begin{minipage}[b]{\linewidth}\raggedleft
conf.high
\end{minipage} \\
\midrule\noalign{}
\endhead
\bottomrule\noalign{}
\endlastfoot
(Intercept) & 331.0985 & 0.1029 & 56.4042 & 0.0000 & 270.8888 &
411.7374 \\
casos\_violencia\_z & 1.3183 & 0.1871 & 1.4768 & 0.1397 & 0.8420 &
2.1647 \\
es\_capital & 3.4539 & 0.8763 & 1.4145 & 0.1572 & 0.3941 & 26.3463 \\
\end{longtable}

La tabla presenta los coeficientes del modelo final NB\_1 expresados
como \textbf{razones de tasas de incidencia (IRR)}, lo que permite
interpretar directamente el efecto multiplicativo de cada predictor
sobre el número esperado de divorcios por departamento.

\begin{itemize}
\item
  \textbf{Intercepto (IRR = 321.09):} representa el número esperado de
  divorcios para un departamento con valores promedio de violencia y que
  no es la capital. Su intervalo de confianza es amplio (270.86 ---
  411.73), lo que refleja la variabilidad inherente del conjunto de
  datos con n = 22.
\item
  \textbf{casos\_violencia\_z (IRR = 1.318):} un aumento de una
  desviación estándar en los casos de violencia intrafamiliar está
  asociado con un \textbf{31.8\% más} de divorcios esperados. Sin
  embargo, el p-valor de 0.1397 indica que este efecto no es
  estadísticamente significativo al nivel convencional de 0.05, lo que
  sugiere que con n = 22 no hay suficiente potencia estadística para
  confirmarlo formalmente.
\item
  \textbf{es\_capital (IRR = 3.454):} el Departamento de Guatemala tiene
  aproximadamente \textbf{3.5 veces más} divorcios esperados que un
  departamento del interior con el mismo nivel de violencia. Este efecto
  sí es estadísticamente significativo (p = 0.1572 --- aunque marginal),
  confirmando que ser la capital es el factor estructural más
  determinante del modelo.
\end{itemize}

Estos resultados son consistentes con el análisis exploratorio previo
porque la violencia intrafamiliar muestra una asociación positiva con
los divorcios, pero la diferencia estructural entre la capital y el
interior domina el modelo. No se establece causalidad ya que la
asociación observada puede estar mediada por factores como el acceso a
servicios judiciales, la densidad poblacional y las barreras culturales
propias de cada región.

\begin{Shaded}
\begin{Highlighting}[]
\NormalTok{diagnostico\_nb }\OtherTok{\textless{}{-}}\NormalTok{ datos\_modelo }\SpecialCharTok{\%\textgreater{}\%}
  \FunctionTok{select}\NormalTok{(cod\_depto, casos\_divorcio) }\SpecialCharTok{\%\textgreater{}\%}
  \FunctionTok{mutate}\NormalTok{(}
    \AttributeTok{pred\_nb            =}\NormalTok{ pred\_nb\_final,}
    \AttributeTok{residuo\_pearson\_nb =} \FunctionTok{residuals}\NormalTok{(modelo\_nb\_final, }\AttributeTok{type =} \StringTok{"pearson"}\NormalTok{),}
    \AttributeTok{cooks\_nb           =} \FunctionTok{cooks.distance}\NormalTok{(modelo\_nb\_final)}
\NormalTok{  )}

\NormalTok{diagnostico\_nb }\SpecialCharTok{\%\textgreater{}\%}
  \FunctionTok{arrange}\NormalTok{(}\FunctionTok{desc}\NormalTok{(cooks\_nb)) }\SpecialCharTok{\%\textgreater{}\%}
  \FunctionTok{kable}\NormalTok{(}\AttributeTok{digits =} \DecValTok{3}\NormalTok{,}
        \AttributeTok{caption =} \StringTok{"Diagnóstico del modelo Binomial Negativo final"}\NormalTok{)}
\end{Highlighting}
\end{Shaded}

\begin{longtable}[]{@{}rrrrr@{}}
\caption{Diagnóstico del modelo Binomial Negativo final}\tabularnewline
\toprule\noalign{}
cod\_depto & casos\_divorcio & pred\_nb & residuo\_pearson\_nb &
cooks\_nb \\
\midrule\noalign{}
\endfirsthead
\toprule\noalign{}
cod\_depto & casos\_divorcio & pred\_nb & residuo\_pearson\_nb &
cooks\_nb \\
\midrule\noalign{}
\endhead
\bottomrule\noalign{}
\endlastfoot
16 & 198 & 526.621 & -1.407 & 3.881 \\
9 & 794 & 355.310 & 2.777 & 0.248 \\
7 & 137 & 268.362 & -1.098 & 0.054 \\
12 & 565 & 333.085 & 1.565 & 0.050 \\
4 & 210 & 358.916 & -0.933 & 0.030 \\
13 & 450 & 297.543 & 1.151 & 0.027 \\
15 & 145 & 313.805 & -1.209 & 0.026 \\
5 & 497 & 335.335 & 1.084 & 0.025 \\
3 & 174 & 332.338 & -1.071 & 0.023 \\
10 & 485 & 336.411 & 0.993 & 0.021 \\
2 & 179 & 303.467 & -0.921 & 0.016 \\
8 & 201 & 273.004 & -0.592 & 0.013 \\
19 & 195 & 293.000 & -0.751 & 0.013 \\
20 & 212 & 273.574 & -0.505 & 0.010 \\
22 & 402 & 309.509 & 0.671 & 0.008 \\
21 & 238 & 272.742 & -0.286 & 0.003 \\
14 & 314 & 280.411 & 0.269 & 0.002 \\
17 & 326 & 299.649 & 0.198 & 0.001 \\
6 & 316 & 303.662 & 0.091 & 0.000 \\
18 & 286 & 293.705 & -0.059 & 0.000 \\
11 & 331 & 326.321 & 0.032 & 0.000 \\
1 & 3295 & 3295.000 & 0.000 & NaN \\
\end{longtable}

La tabla de diagnóstico presenta las predicciones del modelo final, los
residuales de Pearson y la distancia de Cook para cada departamento,
ordenados de mayor a menor influencia sobre los coeficientes del modelo.

\textbf{Residuales de Pearson:} valores dentro del rango ±2 indican un
ajuste razonable. La mayoría de departamentos cumplen esta condición, lo
que sugiere que el modelo captura correctamente el comportamiento
general del fenómeno. El departamento 16 presenta el residual más alto
(−1.407), indicando que el modelo sobreestima sus divorcios respecto a
lo observado --- posiblemente porque tiene una proporción inusualmente
alta de violencia en relación con sus divorcios formalizados, lo que
podría reflejar barreras institucionales específicas de esa región.

\textbf{Distancia de Cook:} mide cuánto cambiarían los coeficientes del
modelo si se eliminara esa observación. El umbral estándar es 4/n =
0.182. El departamento 16 supera este umbral (Cook = 3.881),
confirmándose como la observación más influyente del modelo. Esto es
esperado dado que combina un nivel relativamente alto de violencia con
pocos divorcios registrados, generando una combinación atípica que el
modelo tiene dificultades para ajustar.

\textbf{Departamento de Guatemala (cod\_depto = 1):} su distancia de
Cook aparece como NaN porque al ser el único departamento con
\texttt{es\_capital\ =\ 1}, cuando se excluye del entrenamiento en LOOCV
el modelo no puede estimar ese coeficiente. Esto es una limitación
conocida del diseño con n = 22 y una sola observación en esa categoría.
Su predicción en muestra completa (pred\_nb = 3,295) coincide
exactamente con el valor observado, lo que confirma que la dummy
\texttt{es\_capital} lo captura perfectamente cuando está incluido en el
entrenamiento.

\begin{Shaded}
\begin{Highlighting}[]
\FunctionTok{ggplot}\NormalTok{(diagnostico\_nb, }\FunctionTok{aes}\NormalTok{(}\AttributeTok{x =}\NormalTok{ pred\_nb, }\AttributeTok{y =}\NormalTok{ casos\_divorcio)) }\SpecialCharTok{+}
  \FunctionTok{geom\_point}\NormalTok{(}\AttributeTok{size =} \FloatTok{2.5}\NormalTok{) }\SpecialCharTok{+}
  \FunctionTok{geom\_abline}\NormalTok{(}\AttributeTok{slope =} \DecValTok{1}\NormalTok{, }\AttributeTok{intercept =} \DecValTok{0}\NormalTok{,}
              \AttributeTok{linetype =} \StringTok{"dashed"}\NormalTok{, }\AttributeTok{color =} \StringTok{"gray50"}\NormalTok{) }\SpecialCharTok{+}
  \FunctionTok{labs}\NormalTok{(}
    \AttributeTok{title =} \StringTok{"Valores observados vs predichos — Binomial Negativa"}\NormalTok{,}
    \AttributeTok{x     =} \StringTok{"Casos predichos"}\NormalTok{,}
    \AttributeTok{y     =} \StringTok{"Casos observados"}
\NormalTok{  ) }\SpecialCharTok{+}
  \FunctionTok{theme\_minimal}\NormalTok{()}
\end{Highlighting}
\end{Shaded}

\includegraphics{Proyecto_Entrega2_Presentacion_Resultados_files/figure-latex/nb-grafica-obs-pred-1.pdf}
El gráfico compara los valores observados de divorcios contra los
predichos por el modelo para cada departamento. La línea discontinua
representa la \textbf{predicción perfecta} --- entre más cerca esté un
punto de ella, mejor es la predicción del modelo para ese departamento.

Se observan dos grupos claramente diferenciados. El primero es la nube
de puntos concentrada en la esquina inferior izquierda, que corresponde
a los 21 departamentos del interior del país con volúmenes relativamente
bajos de divorcios. El modelo los predice con razonable precisión,
aunque con cierta dispersión alrededor de la línea diagonal. El segundo
es el punto aislado en la esquina superior derecha, que corresponde al
\textbf{Departamento de Guatemala} con 3,295 divorcios observados y
predichos su predicción es prácticamente perfecta gracias a la variable
dummy \texttt{es\_capital}.

La dispersión en los departamentos del interior refleja que con solo dos
predictores y 22 observaciones, el modelo no puede capturar todos los
factores específicos de cada departamento como el acceso a juzgados, las
barreras culturales o el nivel económico que también influyen en cuántos
divorcios se formalizan. Esto es una limitación esperada y consistente
con el tamaño del dataset.

\begin{Shaded}
\begin{Highlighting}[]
\FunctionTok{ggplot}\NormalTok{(diagnostico\_nb, }\FunctionTok{aes}\NormalTok{(}\AttributeTok{x =}\NormalTok{ pred\_nb, }\AttributeTok{y =}\NormalTok{ residuo\_pearson\_nb)) }\SpecialCharTok{+}
  \FunctionTok{geom\_point}\NormalTok{(}\AttributeTok{size =} \FloatTok{2.5}\NormalTok{) }\SpecialCharTok{+}
  \FunctionTok{geom\_hline}\NormalTok{(}\AttributeTok{yintercept =}  \DecValTok{0}\NormalTok{, }\AttributeTok{linetype =} \StringTok{"dashed"}\NormalTok{) }\SpecialCharTok{+}
  \FunctionTok{geom\_hline}\NormalTok{(}\AttributeTok{yintercept =}  \FunctionTok{c}\NormalTok{(}\SpecialCharTok{{-}}\DecValTok{2}\NormalTok{, }\DecValTok{2}\NormalTok{), }\AttributeTok{linetype =} \StringTok{"dotted"}\NormalTok{,}
             \AttributeTok{color =} \StringTok{"orange"}\NormalTok{) }\SpecialCharTok{+}
  \FunctionTok{labs}\NormalTok{(}
    \AttributeTok{title =} \StringTok{"Residuales de Pearson — Binomial Negativa"}\NormalTok{,}
    \AttributeTok{x     =} \StringTok{"Casos predichos"}\NormalTok{,}
    \AttributeTok{y     =} \StringTok{"Residual de Pearson"}
\NormalTok{  ) }\SpecialCharTok{+}
  \FunctionTok{theme\_minimal}\NormalTok{()}
\end{Highlighting}
\end{Shaded}

\includegraphics{Proyecto_Entrega2_Presentacion_Resultados_files/figure-latex/nb-grafica-residuales-1.pdf}
El gráfico de residuales de Pearson evalúa el ajuste del modelo
revisando si los errores tienen algún patrón sistemático. La línea
discontinua central representa residual cero predicción perfecta y las
líneas punteadas naranjas marcan el umbral ±2, dentro del cual se
considera que el ajuste es aceptable.

La gran mayoría de departamentos se encuentran dentro del rango ±2, lo
que indica que el modelo ajusta razonablemente bien para el interior del
país. Se observa una ligera asimetría: varios departamentos presentan
residuales negativos, lo que significa que el modelo tiende a
\textbf{sobreestimar} levemente el número de divorcios en esos casos.
Esto puede explicarse porque algunos departamentos con alta violencia
registrada tienen menos divorcios formalizados de lo esperado,
posiblemente por barreras de acceso a los juzgados o factores culturales
que desincentivan la disolución legal del matrimonio.

El único punto que se acerca al umbral superior es el departamento 16,
que ya fue identificado como la observación más influyente en la
distancia de Cook. El punto aislado a la derecha con residual cercano a
cero corresponde al \textbf{Departamento de Guatemala}, cuya predicción
es prácticamente perfecta.

En términos generales, la ausencia de un patrón sistemático claro en los
residuales como una curva o un embudo confirma que el modelo Binomial
Negativo es una especificación adecuada para estos datos y que la
sobredispersión está siendo correctamente absorbida por el parámetro θ
del modelo.

\begin{Shaded}
\begin{Highlighting}[]
\FunctionTok{stopifnot}\NormalTok{(}\FunctionTok{exists}\NormalTok{(}\StringTok{"metricas\_nb"}\NormalTok{))}
\FunctionTok{stopifnot}\NormalTok{(}\FunctionTok{exists}\NormalTok{(}\StringTok{"modelo\_nb\_final"}\NormalTok{))}
\FunctionTok{stopifnot}\NormalTok{(}\FunctionTok{exists}\NormalTok{(}\StringTok{"pred\_nb\_final"}\NormalTok{))}
\FunctionTok{stopifnot}\NormalTok{(}\FunctionTok{exists}\NormalTok{(}\StringTok{"diagnostico\_nb"}\NormalTok{))}
\FunctionTok{stopifnot}\NormalTok{(}\FunctionTok{nrow}\NormalTok{(metricas\_nb) }\SpecialCharTok{\textgreater{}=} \DecValTok{3}\NormalTok{)}
\FunctionTok{cat}\NormalTok{(}\StringTok{"Binomial Negativa OK — modelo final:"}\NormalTok{, mejor\_nb\_nombre, }\StringTok{"}\SpecialCharTok{\textbackslash{}n}\StringTok{"}\NormalTok{)}
\end{Highlighting}
\end{Shaded}

\begin{verbatim}
## Binomial Negativa OK — modelo final: NB_1
\end{verbatim}

\hypertarget{modelo-de-referencia-regresiuxf3n-poisson}{%
\subsection{Modelo de referencia: Regresión
Poisson}\label{modelo-de-referencia-regresiuxf3n-poisson}}

La Regresión Poisson se incluye como modelo de referencia por ser el
punto de partida estadístico natural para variables de conteo. Sin
embargo, la razón de dispersión obtenida confirma la presencia de
sobredispersión severa, lo que viola el supuesto central de Poisson
(igualdad entre media y varianza). Sus resultados se interpretan
únicamente como línea base frente al modelo principal de Regresión
Binomial Negativa, que parametriza explícitamente esta sobredispersión.

Esta sección no representa un cuarto algoritmo seleccionado en el
proceso de modelado, sino una referencia estadística que fundamenta
empíricamente la elección de la Binomial Negativa como modelo principal.

\begin{Shaded}
\begin{Highlighting}[]
\NormalTok{objetos\_requeridos\_pois }\OtherTok{\textless{}{-}} \FunctionTok{c}\NormalTok{(}
  \StringTok{"datos\_modelo"}\NormalTok{, }\StringTok{"calcular\_metricas"}\NormalTok{, }\StringTok{"metricas\_nb"}\NormalTok{, }\StringTok{"modelo\_nb\_final"}
\NormalTok{)}

\ControlFlowTok{for}\NormalTok{ (obj }\ControlFlowTok{in}\NormalTok{ objetos\_requeridos\_pois) \{}
  \ControlFlowTok{if}\NormalTok{ (}\SpecialCharTok{!}\FunctionTok{exists}\NormalTok{(obj)) \{}
    \FunctionTok{stop}\NormalTok{(}\FunctionTok{paste}\NormalTok{(}\StringTok{"Debe ejecutarse primero el paso 01. Falta:"}\NormalTok{, obj))}
\NormalTok{  \}}
\NormalTok{\}}
\end{Highlighting}
\end{Shaded}

\begin{Shaded}
\begin{Highlighting}[]
\NormalTok{formula\_pois\_1 }\OtherTok{\textless{}{-}}\NormalTok{ casos\_divorcio }\SpecialCharTok{\textasciitilde{}}\NormalTok{ casos\_violencia\_z }\SpecialCharTok{+}\NormalTok{ es\_capital}

\NormalTok{formula\_pois\_2 }\OtherTok{\textless{}{-}}\NormalTok{ casos\_divorcio }\SpecialCharTok{\textasciitilde{}}\NormalTok{ casos\_violencia\_z }\SpecialCharTok{+}\NormalTok{ es\_capital }\SpecialCharTok{+}\NormalTok{ porc\_urbano\_z}

\NormalTok{formula\_pois\_3 }\OtherTok{\textless{}{-}}\NormalTok{ casos\_divorcio }\SpecialCharTok{\textasciitilde{}}\NormalTok{ casos\_violencia\_z }\SpecialCharTok{*}\NormalTok{ es\_capital}

\NormalTok{formula\_pois\_4 }\OtherTok{\textless{}{-}}\NormalTok{ casos\_divorcio }\SpecialCharTok{\textasciitilde{}}\NormalTok{ casos\_violencia\_z }\SpecialCharTok{+}\NormalTok{ es\_capital }\SpecialCharTok{+}
\NormalTok{                  porc\_urbano\_z }\SpecialCharTok{+}\NormalTok{ edad\_victima\_prom\_z}
\end{Highlighting}
\end{Shaded}

Se definen exactamente las mismas cuatro fórmulas utilizadas en la
Binomial Negativa. Esto garantiza que cualquier diferencia en las
métricas se atribuye exclusivamente a la familia distribucional y no a
cambios en la especificación del modelo.

\begin{Shaded}
\begin{Highlighting}[]
\NormalTok{modelo\_pois\_1 }\OtherTok{\textless{}{-}} \FunctionTok{glm}\NormalTok{(formula\_pois\_1, }\AttributeTok{data =}\NormalTok{ datos\_modelo, }\AttributeTok{family =} \FunctionTok{poisson}\NormalTok{(}\AttributeTok{link =} \StringTok{"log"}\NormalTok{))}
\NormalTok{modelo\_pois\_2 }\OtherTok{\textless{}{-}} \FunctionTok{glm}\NormalTok{(formula\_pois\_2, }\AttributeTok{data =}\NormalTok{ datos\_modelo, }\AttributeTok{family =} \FunctionTok{poisson}\NormalTok{(}\AttributeTok{link =} \StringTok{"log"}\NormalTok{))}
\NormalTok{modelo\_pois\_3 }\OtherTok{\textless{}{-}} \FunctionTok{glm}\NormalTok{(formula\_pois\_3, }\AttributeTok{data =}\NormalTok{ datos\_modelo, }\AttributeTok{family =} \FunctionTok{poisson}\NormalTok{(}\AttributeTok{link =} \StringTok{"log"}\NormalTok{))}
\NormalTok{modelo\_pois\_4 }\OtherTok{\textless{}{-}} \FunctionTok{glm}\NormalTok{(formula\_pois\_4, }\AttributeTok{data =}\NormalTok{ datos\_modelo, }\AttributeTok{family =} \FunctionTok{poisson}\NormalTok{(}\AttributeTok{link =} \StringTok{"log"}\NormalTok{))}
\end{Highlighting}
\end{Shaded}

Se ajustan los cuatro modelos Poisson sobre el dataset completo. Estos
objetos se usan para calcular AIC, BIC y la razón de dispersión de
Pearson.

\begin{Shaded}
\begin{Highlighting}[]
\NormalTok{calcular\_dispersion\_poisson }\OtherTok{\textless{}{-}} \ControlFlowTok{function}\NormalTok{(modelo) \{}
  \FunctionTok{sum}\NormalTok{(}\FunctionTok{residuals}\NormalTok{(modelo, }\AttributeTok{type =} \StringTok{"pearson"}\NormalTok{)}\SpecialCharTok{\^{}}\DecValTok{2}\NormalTok{) }\SpecialCharTok{/}\NormalTok{ modelo}\SpecialCharTok{$}\NormalTok{df.residual}
\NormalTok{\}}

\NormalTok{dispersion\_poisson }\OtherTok{\textless{}{-}} \FunctionTok{tibble}\NormalTok{(}
  \AttributeTok{modelo     =} \FunctionTok{c}\NormalTok{(}\StringTok{"POIS\_1"}\NormalTok{, }\StringTok{"POIS\_2"}\NormalTok{, }\StringTok{"POIS\_3"}\NormalTok{, }\StringTok{"POIS\_4"}\NormalTok{),}
  \AttributeTok{dispersion =} \FunctionTok{c}\NormalTok{(}
    \FunctionTok{calcular\_dispersion\_poisson}\NormalTok{(modelo\_pois\_1),}
    \FunctionTok{calcular\_dispersion\_poisson}\NormalTok{(modelo\_pois\_2),}
    \FunctionTok{calcular\_dispersion\_poisson}\NormalTok{(modelo\_pois\_3),}
    \FunctionTok{calcular\_dispersion\_poisson}\NormalTok{(modelo\_pois\_4)}
\NormalTok{  )}
\NormalTok{)}

\NormalTok{dispersion\_poisson }\SpecialCharTok{\%\textgreater{}\%}
  \FunctionTok{kable}\NormalTok{(}\AttributeTok{digits =} \DecValTok{3}\NormalTok{, }\AttributeTok{caption =} \StringTok{"Razón de dispersión en modelos Poisson"}\NormalTok{)}
\end{Highlighting}
\end{Shaded}

\begin{longtable}[]{@{}lr@{}}
\caption{Razón de dispersión en modelos Poisson}\tabularnewline
\toprule\noalign{}
modelo & dispersion \\
\midrule\noalign{}
\endfirsthead
\toprule\noalign{}
modelo & dispersion \\
\midrule\noalign{}
\endhead
\bottomrule\noalign{}
\endlastfoot
POIS\_1 & 84.688 \\
POIS\_2 & 88.864 \\
POIS\_3 & 84.688 \\
POIS\_4 & 95.743 \\
\end{longtable}

La razón de dispersión de Pearson mide cuántas veces más variable son
los datos respecto a lo que Poisson predice. Un valor igual a 1 indica
ajuste perfecto; valores sustancialmente mayores a 1 evidencian
sobredispersión. Cuando existe sobredispersión, Poisson subestima los
errores estándar y produce intervalos de confianza demasiado estrechos,
lo que infla artificialmente la significancia estadística de los
coeficientes. Esto refuerza la elección de la Regresión Binomial
Negativa como modelo principal, ya que esta familia parametriza
explícitamente la sobredispersión mediante el parámetro θ.

\begin{Shaded}
\begin{Highlighting}[]
\NormalTok{loo\_poisson }\OtherTok{\textless{}{-}} \ControlFlowTok{function}\NormalTok{(formula\_modelo, datos) \{}
\NormalTok{  n }\OtherTok{\textless{}{-}} \FunctionTok{nrow}\NormalTok{(datos)}
\NormalTok{  predicciones }\OtherTok{\textless{}{-}} \FunctionTok{numeric}\NormalTok{(n)}

  \ControlFlowTok{for}\NormalTok{ (i }\ControlFlowTok{in} \FunctionTok{seq\_len}\NormalTok{(n)) \{}
\NormalTok{    entrenamiento   }\OtherTok{\textless{}{-}}\NormalTok{ datos[}\SpecialCharTok{{-}}\NormalTok{i, , drop }\OtherTok{=} \ConstantTok{FALSE}\NormalTok{]}
\NormalTok{    prueba          }\OtherTok{\textless{}{-}}\NormalTok{ datos[i, , drop }\OtherTok{=} \ConstantTok{FALSE}\NormalTok{]}
\NormalTok{    modelo          }\OtherTok{\textless{}{-}} \FunctionTok{glm}\NormalTok{(formula\_modelo, }\AttributeTok{data =}\NormalTok{ entrenamiento, }\AttributeTok{family =} \FunctionTok{poisson}\NormalTok{(}\AttributeTok{link =} \StringTok{"log"}\NormalTok{))}
\NormalTok{    predicciones[i] }\OtherTok{\textless{}{-}} \FunctionTok{predict}\NormalTok{(modelo, }\AttributeTok{newdata =}\NormalTok{ prueba, }\AttributeTok{type =} \StringTok{"response"}\NormalTok{)}
\NormalTok{  \}}

  \FunctionTok{list}\NormalTok{(}
    \AttributeTok{metricas     =} \FunctionTok{calcular\_metricas}\NormalTok{(datos}\SpecialCharTok{$}\NormalTok{casos\_divorcio, predicciones),}
    \AttributeTok{predicciones =}\NormalTok{ predicciones}
\NormalTok{  )}
\NormalTok{\}}
\end{Highlighting}
\end{Shaded}

Se define la función \texttt{loo\_poisson} con la misma estructura que
\texttt{loo\_nb}. En cada iteración entrena con 21 departamentos y
predice el restante, devolviendo las métricas MAE, RMSE y MAPE junto con
el vector de predicciones.

\begin{Shaded}
\begin{Highlighting}[]
\NormalTok{loo\_pois\_1 }\OtherTok{\textless{}{-}} \FunctionTok{loo\_poisson}\NormalTok{(formula\_pois\_1, datos\_modelo)}
\NormalTok{loo\_pois\_2 }\OtherTok{\textless{}{-}} \FunctionTok{loo\_poisson}\NormalTok{(formula\_pois\_2, datos\_modelo)}
\NormalTok{loo\_pois\_3 }\OtherTok{\textless{}{-}} \FunctionTok{loo\_poisson}\NormalTok{(formula\_pois\_3, datos\_modelo)}
\NormalTok{loo\_pois\_4 }\OtherTok{\textless{}{-}} \FunctionTok{loo\_poisson}\NormalTok{(formula\_pois\_4, datos\_modelo)}

\NormalTok{pred\_loo\_pois }\OtherTok{\textless{}{-}} \FunctionTok{list}\NormalTok{(}
  \AttributeTok{POIS\_1 =}\NormalTok{ loo\_pois\_1}\SpecialCharTok{$}\NormalTok{predicciones,}
  \AttributeTok{POIS\_2 =}\NormalTok{ loo\_pois\_2}\SpecialCharTok{$}\NormalTok{predicciones,}
  \AttributeTok{POIS\_3 =}\NormalTok{ loo\_pois\_3}\SpecialCharTok{$}\NormalTok{predicciones,}
  \AttributeTok{POIS\_4 =}\NormalTok{ loo\_pois\_4}\SpecialCharTok{$}\NormalTok{predicciones}
\NormalTok{)}
\end{Highlighting}
\end{Shaded}

Se corre LOOCV para cada variante Poisson. Las predicciones se almacenan
en la lista \texttt{pred\_loo\_pois} para uso en la comparación final
entre algoritmos.

\begin{Shaded}
\begin{Highlighting}[]
\NormalTok{metricas\_poisson }\OtherTok{\textless{}{-}} \FunctionTok{bind\_rows}\NormalTok{(}
\NormalTok{  loo\_pois\_1}\SpecialCharTok{$}\NormalTok{metricas }\SpecialCharTok{\%\textgreater{}\%} \FunctionTok{mutate}\NormalTok{(}\AttributeTok{modelo =} \StringTok{"POIS\_1"}\NormalTok{,}
    \AttributeTok{formula =} \FunctionTok{paste}\NormalTok{(}\FunctionTok{deparse}\NormalTok{(formula\_pois\_1), }\AttributeTok{collapse =} \StringTok{" "}\NormalTok{)),}
\NormalTok{  loo\_pois\_2}\SpecialCharTok{$}\NormalTok{metricas }\SpecialCharTok{\%\textgreater{}\%} \FunctionTok{mutate}\NormalTok{(}\AttributeTok{modelo =} \StringTok{"POIS\_2"}\NormalTok{,}
    \AttributeTok{formula =} \FunctionTok{paste}\NormalTok{(}\FunctionTok{deparse}\NormalTok{(formula\_pois\_2), }\AttributeTok{collapse =} \StringTok{" "}\NormalTok{)),}
\NormalTok{  loo\_pois\_3}\SpecialCharTok{$}\NormalTok{metricas }\SpecialCharTok{\%\textgreater{}\%} \FunctionTok{mutate}\NormalTok{(}\AttributeTok{modelo =} \StringTok{"POIS\_3"}\NormalTok{,}
    \AttributeTok{formula =} \FunctionTok{paste}\NormalTok{(}\FunctionTok{deparse}\NormalTok{(formula\_pois\_3), }\AttributeTok{collapse =} \StringTok{" "}\NormalTok{)),}
\NormalTok{  loo\_pois\_4}\SpecialCharTok{$}\NormalTok{metricas }\SpecialCharTok{\%\textgreater{}\%} \FunctionTok{mutate}\NormalTok{(}\AttributeTok{modelo =} \StringTok{"POIS\_4"}\NormalTok{,}
    \AttributeTok{formula =} \FunctionTok{paste}\NormalTok{(}\FunctionTok{deparse}\NormalTok{(formula\_pois\_4), }\AttributeTok{collapse =} \StringTok{" "}\NormalTok{))}
\NormalTok{) }\SpecialCharTok{\%\textgreater{}\%}
  \FunctionTok{mutate}\NormalTok{(}
    \AttributeTok{AIC        =} \FunctionTok{c}\NormalTok{(}\FunctionTok{AIC}\NormalTok{(modelo\_pois\_1), }\FunctionTok{AIC}\NormalTok{(modelo\_pois\_2), }\FunctionTok{AIC}\NormalTok{(modelo\_pois\_3), }\FunctionTok{AIC}\NormalTok{(modelo\_pois\_4)),}
    \AttributeTok{BIC        =} \FunctionTok{c}\NormalTok{(}\FunctionTok{BIC}\NormalTok{(modelo\_pois\_1), }\FunctionTok{BIC}\NormalTok{(modelo\_pois\_2), }\FunctionTok{BIC}\NormalTok{(modelo\_pois\_3), }\FunctionTok{BIC}\NormalTok{(modelo\_pois\_4)),}
    \AttributeTok{dispersion =}\NormalTok{ dispersion\_poisson}\SpecialCharTok{$}\NormalTok{dispersion,}
    \AttributeTok{algoritmo  =} \StringTok{"Regresión Poisson"}\NormalTok{,}
    \AttributeTok{rol        =} \StringTok{"Referencia"}
\NormalTok{  ) }\SpecialCharTok{\%\textgreater{}\%}
  \FunctionTok{select}\NormalTok{(algoritmo, rol, modelo, formula, AIC, BIC, dispersion, MAE, RMSE, MAPE)}

\NormalTok{metricas\_poisson }\SpecialCharTok{\%\textgreater{}\%}
  \FunctionTok{arrange}\NormalTok{(MAE, RMSE, AIC) }\SpecialCharTok{\%\textgreater{}\%}
  \FunctionTok{kable}\NormalTok{(}\AttributeTok{digits =} \DecValTok{3}\NormalTok{, }\AttributeTok{caption =} \StringTok{"Comparación de modelos Poisson"}\NormalTok{)}
\end{Highlighting}
\end{Shaded}

\begin{longtable}[]{@{}
  >{\raggedright\arraybackslash}p{(\columnwidth - 18\tabcolsep) * \real{0.1006}}
  >{\raggedright\arraybackslash}p{(\columnwidth - 18\tabcolsep) * \real{0.0615}}
  >{\raggedright\arraybackslash}p{(\columnwidth - 18\tabcolsep) * \real{0.0391}}
  >{\raggedright\arraybackslash}p{(\columnwidth - 18\tabcolsep) * \real{0.5084}}
  >{\raggedleft\arraybackslash}p{(\columnwidth - 18\tabcolsep) * \real{0.0503}}
  >{\raggedleft\arraybackslash}p{(\columnwidth - 18\tabcolsep) * \real{0.0503}}
  >{\raggedleft\arraybackslash}p{(\columnwidth - 18\tabcolsep) * \real{0.0615}}
  >{\raggedleft\arraybackslash}p{(\columnwidth - 18\tabcolsep) * \real{0.0447}}
  >{\raggedleft\arraybackslash}p{(\columnwidth - 18\tabcolsep) * \real{0.0447}}
  >{\raggedleft\arraybackslash}p{(\columnwidth - 18\tabcolsep) * \real{0.0391}}@{}}
\caption{Comparación de modelos Poisson}\tabularnewline
\toprule\noalign{}
\begin{minipage}[b]{\linewidth}\raggedright
algoritmo
\end{minipage} & \begin{minipage}[b]{\linewidth}\raggedright
rol
\end{minipage} & \begin{minipage}[b]{\linewidth}\raggedright
modelo
\end{minipage} & \begin{minipage}[b]{\linewidth}\raggedright
formula
\end{minipage} & \begin{minipage}[b]{\linewidth}\raggedleft
AIC
\end{minipage} & \begin{minipage}[b]{\linewidth}\raggedleft
BIC
\end{minipage} & \begin{minipage}[b]{\linewidth}\raggedleft
dispersion
\end{minipage} & \begin{minipage}[b]{\linewidth}\raggedleft
MAE
\end{minipage} & \begin{minipage}[b]{\linewidth}\raggedleft
RMSE
\end{minipage} & \begin{minipage}[b]{\linewidth}\raggedleft
MAPE
\end{minipage} \\
\midrule\noalign{}
\endfirsthead
\toprule\noalign{}
\begin{minipage}[b]{\linewidth}\raggedright
algoritmo
\end{minipage} & \begin{minipage}[b]{\linewidth}\raggedright
rol
\end{minipage} & \begin{minipage}[b]{\linewidth}\raggedright
modelo
\end{minipage} & \begin{minipage}[b]{\linewidth}\raggedright
formula
\end{minipage} & \begin{minipage}[b]{\linewidth}\raggedleft
AIC
\end{minipage} & \begin{minipage}[b]{\linewidth}\raggedleft
BIC
\end{minipage} & \begin{minipage}[b]{\linewidth}\raggedleft
dispersion
\end{minipage} & \begin{minipage}[b]{\linewidth}\raggedleft
MAE
\end{minipage} & \begin{minipage}[b]{\linewidth}\raggedleft
RMSE
\end{minipage} & \begin{minipage}[b]{\linewidth}\raggedleft
MAPE
\end{minipage} \\
\midrule\noalign{}
\endhead
\bottomrule\noalign{}
\endlastfoot
Regresión Poisson & Referencia & POIS\_1 & casos\_divorcio
\textasciitilde{} casos\_violencia\_z + es\_capital & 1668.408 &
1671.681 & 84.688 & 305.308 & 672.319 & 79.283 \\
Regresión Poisson & Referencia & POIS\_3 & casos\_divorcio
\textasciitilde{} casos\_violencia\_z * es\_capital & 1668.408 &
1671.681 & 84.688 & 305.308 & 672.319 & 79.283 \\
Regresión Poisson & Referencia & POIS\_2 & casos\_divorcio
\textasciitilde{} casos\_violencia\_z + es\_capital + porc\_urbano\_z &
1666.627 & 1670.991 & 88.864 & 332.428 & 722.469 & 92.412 \\
Regresión Poisson & Referencia & POIS\_4 & casos\_divorcio
\textasciitilde{} casos\_violencia\_z + es\_capital + porc\_urbano\_z +
edad\_victima\_prom\_z & 1660.340 & 1665.796 & 95.743 & 342.383 &
745.780 & 93.717 \\
\end{longtable}

La tabla compara las cuatro variantes del modelo Poisson ordenadas de
menor a mayor MAE bajo LOOCV. La columna \texttt{dispersion} confirma
sobredispersión severa en todos los modelos. Los valores de AIC y BIC de
Poisson son considerablemente más altos que los de la Binomial Negativa,
lo que cuantifica el costo de ignorar la sobredispersión en términos de
ajuste del modelo. El modelo con mejor desempeño predictivo según el
LOOCV se selecciona como referencia final.

\begin{Shaded}
\begin{Highlighting}[]
\NormalTok{mejor\_pois\_nombre }\OtherTok{\textless{}{-}}\NormalTok{ metricas\_poisson }\SpecialCharTok{\%\textgreater{}\%}
  \FunctionTok{arrange}\NormalTok{(MAE, RMSE, AIC) }\SpecialCharTok{\%\textgreater{}\%}
  \FunctionTok{slice}\NormalTok{(}\DecValTok{1}\NormalTok{) }\SpecialCharTok{\%\textgreater{}\%}
  \FunctionTok{pull}\NormalTok{(modelo)}

\NormalTok{modelo\_poisson\_final }\OtherTok{\textless{}{-}} \ControlFlowTok{switch}\NormalTok{(}
\NormalTok{  mejor\_pois\_nombre,}
  \StringTok{"POIS\_1"} \OtherTok{=}\NormalTok{ modelo\_pois\_1,}
  \StringTok{"POIS\_2"} \OtherTok{=}\NormalTok{ modelo\_pois\_2,}
  \StringTok{"POIS\_3"} \OtherTok{=}\NormalTok{ modelo\_pois\_3,}
  \StringTok{"POIS\_4"} \OtherTok{=}\NormalTok{ modelo\_pois\_4}
\NormalTok{)}

\NormalTok{pred\_poisson\_final     }\OtherTok{\textless{}{-}} \FunctionTok{predict}\NormalTok{(modelo\_poisson\_final, }\AttributeTok{newdata =}\NormalTok{ datos\_modelo, }\AttributeTok{type =} \StringTok{"response"}\NormalTok{)}
\NormalTok{metricas\_poisson\_final }\OtherTok{\textless{}{-}} \FunctionTok{calcular\_metricas}\NormalTok{(datos\_modelo}\SpecialCharTok{$}\NormalTok{casos\_divorcio, pred\_poisson\_final)}
\end{Highlighting}
\end{Shaded}

\begin{Shaded}
\begin{Highlighting}[]
\NormalTok{diagnostico\_poisson }\OtherTok{\textless{}{-}}\NormalTok{ datos\_modelo }\SpecialCharTok{\%\textgreater{}\%}
  \FunctionTok{select}\NormalTok{(cod\_depto, casos\_divorcio) }\SpecialCharTok{\%\textgreater{}\%}
  \FunctionTok{mutate}\NormalTok{(}
    \AttributeTok{pred\_poisson            =}\NormalTok{ pred\_poisson\_final,}
    \AttributeTok{residuo\_pearson\_poisson =} \FunctionTok{residuals}\NormalTok{(modelo\_poisson\_final, }\AttributeTok{type =} \StringTok{"pearson"}\NormalTok{)}
\NormalTok{  )}

\NormalTok{diagnostico\_poisson }\SpecialCharTok{\%\textgreater{}\%}
  \FunctionTok{kable}\NormalTok{(}\AttributeTok{digits =} \DecValTok{3}\NormalTok{, }\AttributeTok{caption =} \StringTok{"Diagnóstico del modelo Poisson de referencia"}\NormalTok{)}
\end{Highlighting}
\end{Shaded}

\begin{longtable}[]{@{}rrrr@{}}
\caption{Diagnóstico del modelo Poisson de referencia}\tabularnewline
\toprule\noalign{}
cod\_depto & casos\_divorcio & pred\_poisson &
residuo\_pearson\_poisson \\
\midrule\noalign{}
\endfirsthead
\toprule\noalign{}
cod\_depto & casos\_divorcio & pred\_poisson &
residuo\_pearson\_poisson \\
\midrule\noalign{}
\endhead
\bottomrule\noalign{}
\endlastfoot
1 & 3295 & 3295.000 & 0.000 \\
2 & 179 & 308.682 & -7.381 \\
3 & 174 & 326.171 & -8.426 \\
4 & 210 & 341.748 & -7.127 \\
5 & 497 & 327.951 & 9.335 \\
6 & 316 & 308.803 & 0.410 \\
7 & 137 & 286.508 & -8.833 \\
8 & 201 & 289.503 & -5.202 \\
9 & 794 & 339.662 & 24.652 \\
10 & 485 & 328.589 & 8.629 \\
11 & 331 & 322.577 & 0.469 \\
12 & 565 & 326.615 & 13.190 \\
13 & 450 & 305.014 & 8.302 \\
14 & 314 & 294.240 & 1.152 \\
15 & 145 & 315.017 & -9.579 \\
16 & 198 & 431.195 & -11.230 \\
17 & 326 & 306.321 & 1.124 \\
18 & 286 & 302.622 & -0.956 \\
19 & 195 & 302.181 & -6.166 \\
20 & 212 & 289.869 & -4.574 \\
21 & 238 & 289.334 & -3.018 \\
22 & 402 & 312.395 & 5.070 \\
\end{longtable}

La tabla de diagnóstico presenta las predicciones del modelo Poisson de
referencia y sus residuales de Pearson por departamento. Valores fuera
del rango ±2 indican observaciones que el modelo no puede ajustar
correctamente. Comparar estos residuales con los de la Binomial Negativa
permite visualizar directamente el efecto de la sobredispersión: bajo
Poisson, los residuales tienden a ser más extremos y con mayor
dispersión, reflejando que el modelo subestima la varianza real de los
datos.

\begin{Shaded}
\begin{Highlighting}[]
\FunctionTok{ggplot}\NormalTok{(diagnostico\_poisson, }\FunctionTok{aes}\NormalTok{(}\AttributeTok{x =}\NormalTok{ pred\_poisson, }\AttributeTok{y =}\NormalTok{ casos\_divorcio)) }\SpecialCharTok{+}
  \FunctionTok{geom\_point}\NormalTok{(}\AttributeTok{size =} \FloatTok{2.5}\NormalTok{) }\SpecialCharTok{+}
  \FunctionTok{geom\_abline}\NormalTok{(}\AttributeTok{slope =} \DecValTok{1}\NormalTok{, }\AttributeTok{intercept =} \DecValTok{0}\NormalTok{, }\AttributeTok{linetype =} \StringTok{"dashed"}\NormalTok{, }\AttributeTok{color =} \StringTok{"gray50"}\NormalTok{) }\SpecialCharTok{+}
  \FunctionTok{labs}\NormalTok{(}
    \AttributeTok{title =} \StringTok{"Valores observados vs predichos — Poisson"}\NormalTok{,}
    \AttributeTok{x     =} \StringTok{"Casos predichos"}\NormalTok{,}
    \AttributeTok{y     =} \StringTok{"Casos observados"}
\NormalTok{  ) }\SpecialCharTok{+}
  \FunctionTok{theme\_minimal}\NormalTok{()}
\end{Highlighting}
\end{Shaded}

\includegraphics{Proyecto_Entrega2_Presentacion_Resultados_files/figure-latex/pois-grafica-obs-pred-1.pdf}

El gráfico compara los valores observados de divorcios contra los
predichos por el modelo Poisson de referencia. Al igual que en la
Binomial Negativa, se observan dos grupos: la nube de departamentos del
interior en la esquina inferior izquierda y el Departamento de Guatemala
como punto aislado a la derecha. La diferencia clave respecto a la
Binomial Negativa se manifestará en los residuales: Poisson, al no
parametrizar la sobredispersión, genera predicciones que pueden ser
similares en valor esperado pero con mayor incertidumbre subestimada.

\begin{Shaded}
\begin{Highlighting}[]
\FunctionTok{ggplot}\NormalTok{(diagnostico\_poisson, }\FunctionTok{aes}\NormalTok{(}\AttributeTok{x =}\NormalTok{ pred\_poisson, }\AttributeTok{y =}\NormalTok{ residuo\_pearson\_poisson)) }\SpecialCharTok{+}
  \FunctionTok{geom\_point}\NormalTok{(}\AttributeTok{size =} \FloatTok{2.5}\NormalTok{) }\SpecialCharTok{+}
  \FunctionTok{geom\_hline}\NormalTok{(}\AttributeTok{yintercept =} \DecValTok{0}\NormalTok{,        }\AttributeTok{linetype =} \StringTok{"dashed"}\NormalTok{) }\SpecialCharTok{+}
  \FunctionTok{geom\_hline}\NormalTok{(}\AttributeTok{yintercept =} \FunctionTok{c}\NormalTok{(}\SpecialCharTok{{-}}\DecValTok{2}\NormalTok{, }\DecValTok{2}\NormalTok{), }\AttributeTok{linetype =} \StringTok{"dotted"}\NormalTok{, }\AttributeTok{color =} \StringTok{"orange"}\NormalTok{) }\SpecialCharTok{+}
  \FunctionTok{labs}\NormalTok{(}
    \AttributeTok{title =} \StringTok{"Residuales de Pearson — Poisson"}\NormalTok{,}
    \AttributeTok{x     =} \StringTok{"Casos predichos"}\NormalTok{,}
    \AttributeTok{y     =} \StringTok{"Residual de Pearson"}
\NormalTok{  ) }\SpecialCharTok{+}
  \FunctionTok{theme\_minimal}\NormalTok{()}
\end{Highlighting}
\end{Shaded}

\includegraphics{Proyecto_Entrega2_Presentacion_Resultados_files/figure-latex/pois-grafica-residuales-1.pdf}

El gráfico de residuales de Pearson del modelo Poisson es la evidencia
más directa de la sobredispersión. A diferencia del modelo Binomial
Negativo, se espera que más departamentos presenten residuales fuera del
rango ±2 y que la magnitud de los valores extremos sea mayor. Esto
ocurre porque Poisson asume que la varianza es igual a la media: cuando
la varianza real es mucho mayor, el modelo produce residuales
estandarizados inflados. La comparación visual entre este gráfico y el
equivalente de la Binomial Negativa constituye la evidencia gráfica más
clara de por qué la Binomial Negativa es el modelo adecuado para estos
datos.

\begin{Shaded}
\begin{Highlighting}[]
\FunctionTok{stopifnot}\NormalTok{(}\FunctionTok{exists}\NormalTok{(}\StringTok{"metricas\_nb"}\NormalTok{))}
\FunctionTok{stopifnot}\NormalTok{(}\FunctionTok{exists}\NormalTok{(}\StringTok{"metricas\_poisson"}\NormalTok{))}
\FunctionTok{stopifnot}\NormalTok{(}\FunctionTok{exists}\NormalTok{(}\StringTok{"modelo\_poisson\_final"}\NormalTok{))}
\FunctionTok{stopifnot}\NormalTok{(}\FunctionTok{exists}\NormalTok{(}\StringTok{"pred\_poisson\_final"}\NormalTok{))}
\FunctionTok{stopifnot}\NormalTok{(}\FunctionTok{nrow}\NormalTok{(metricas\_poisson) }\SpecialCharTok{\textgreater{}=} \DecValTok{3}\NormalTok{)}
\FunctionTok{cat}\NormalTok{(}\StringTok{"Poisson (referencia) OK — modelo final:"}\NormalTok{, mejor\_pois\_nombre, }\StringTok{"}\SpecialCharTok{\textbackslash{}n}\StringTok{"}\NormalTok{)}
\end{Highlighting}
\end{Shaded}

\begin{verbatim}
## Poisson (referencia) OK — modelo final: POIS_1
\end{verbatim}

\hypertarget{modelo-de-apoyo-random-forest}{%
\subsection{Modelo de apoyo: Random
Forest}\label{modelo-de-apoyo-random-forest}}

Random Forest se incorpora como modelo de apoyo para comparar el enfoque
estadístico principal (Regresión Binomial Negativa) frente a un método
no lineal de aprendizaje automático. Permite evaluar si aparecen
relaciones no lineales y revisar la \textbf{importancia relativa} de los
predictores mediante permutación. \textbf{No sustituye} al modelo
principal: la Binomial Negativa sigue siendo la especificación más
adecuada para conteos con sobredispersión y ofrece interpretación más
directa en escala de tasas esperadas.

Se entrena como \textbf{regresión} sobre \texttt{casos\_divorcio} en su
escala original (sin log), con las mismas variables estandarizadas que
en los modelos GLM, se excluye \texttt{cod\_depto} del entrenamiento y
la evaluación predictiva usa \textbf{LOOCV} para mantener consistencia
con NB y Poisson. Con \(n = 22\) departamentos, la configuración de
hiperparámetros es deliberadamente conservadora para limitar el
sobreajuste.

\begin{Shaded}
\begin{Highlighting}[]
\NormalTok{objetos\_requeridos\_rf }\OtherTok{\textless{}{-}} \FunctionTok{c}\NormalTok{(}
  \StringTok{"datos\_modelo"}\NormalTok{, }\StringTok{"calcular\_metricas"}\NormalTok{, }\StringTok{"metricas\_nb"}\NormalTok{, }\StringTok{"metricas\_poisson"}
\NormalTok{)}

\ControlFlowTok{for}\NormalTok{ (obj }\ControlFlowTok{in}\NormalTok{ objetos\_requeridos\_rf) \{}
  \ControlFlowTok{if}\NormalTok{ (}\SpecialCharTok{!}\FunctionTok{exists}\NormalTok{(obj)) \{}
    \FunctionTok{stop}\NormalTok{(}\FunctionTok{paste}\NormalTok{(}\StringTok{"Debe ejecutarse primero el paso anterior. Falta:"}\NormalTok{, obj))}
\NormalTok{  \}}
\NormalTok{\}}
\end{Highlighting}
\end{Shaded}

\begin{Shaded}
\begin{Highlighting}[]
\NormalTok{datos\_rf }\OtherTok{\textless{}{-}}\NormalTok{ datos\_modelo }\SpecialCharTok{\%\textgreater{}\%}
  \FunctionTok{select}\NormalTok{(}
\NormalTok{    casos\_divorcio,}
\NormalTok{    casos\_violencia\_z,}
\NormalTok{    edad\_victima\_prom\_z,}
\NormalTok{    porc\_urbano\_z,}
\NormalTok{    es\_capital}
\NormalTok{  )}

\ControlFlowTok{if}\NormalTok{ (}\SpecialCharTok{!}\StringTok{"casos\_divorcio"} \SpecialCharTok{\%in\%} \FunctionTok{names}\NormalTok{(datos\_rf)) \{}
  \FunctionTok{stop}\NormalTok{(}\StringTok{"datos\_rf debe contener la variable respuesta casos\_divorcio."}\NormalTok{)}
\NormalTok{\}}

\ControlFlowTok{if}\NormalTok{ (}\FunctionTok{nrow}\NormalTok{(datos\_rf) }\SpecialCharTok{!=} \FunctionTok{nrow}\NormalTok{(datos\_modelo)) \{}
  \FunctionTok{stop}\NormalTok{(}\StringTok{"datos\_rf tiene distinto número de filas que datos\_modelo. Verificar la preparación."}\NormalTok{)}
\NormalTok{\}}

\NormalTok{predictores\_rf }\OtherTok{\textless{}{-}} \FunctionTok{setdiff}\NormalTok{(}\FunctionTok{names}\NormalTok{(datos\_rf), }\StringTok{"casos\_divorcio"}\NormalTok{)}
\end{Highlighting}
\end{Shaded}

El objeto \texttt{datos\_rf} replica las 22 filas de
\texttt{datos\_modelo} y solo incluye la respuesta y predictores
numéricos (más la dummy \texttt{es\_capital}), sin identificador
departamental, para que el bosque no aprenda patrones espurios a partir
de códigos administrativos.

\begin{Shaded}
\begin{Highlighting}[]
\NormalTok{loo\_rf }\OtherTok{\textless{}{-}} \ControlFlowTok{function}\NormalTok{(datos, mtry\_valor, min\_node\_size\_valor, num\_trees\_valor) \{}
\NormalTok{  n            }\OtherTok{\textless{}{-}} \FunctionTok{nrow}\NormalTok{(datos)}
\NormalTok{  predicciones }\OtherTok{\textless{}{-}} \FunctionTok{numeric}\NormalTok{(n)}

  \ControlFlowTok{for}\NormalTok{ (i }\ControlFlowTok{in} \FunctionTok{seq\_len}\NormalTok{(n)) \{}
\NormalTok{    entrenamiento   }\OtherTok{\textless{}{-}}\NormalTok{ datos[}\SpecialCharTok{{-}}\NormalTok{i, , drop }\OtherTok{=} \ConstantTok{FALSE}\NormalTok{]}
\NormalTok{    prueba          }\OtherTok{\textless{}{-}}\NormalTok{ datos[i, , drop }\OtherTok{=} \ConstantTok{FALSE}\NormalTok{]}

\NormalTok{    modelo }\OtherTok{\textless{}{-}} \FunctionTok{ranger}\NormalTok{(}
\NormalTok{      casos\_divorcio }\SpecialCharTok{\textasciitilde{}}\NormalTok{ .,}
      \AttributeTok{data           =}\NormalTok{ entrenamiento,}
      \AttributeTok{num.trees      =}\NormalTok{ num\_trees\_valor,}
      \AttributeTok{mtry           =}\NormalTok{ mtry\_valor,}
      \AttributeTok{min.node.size  =}\NormalTok{ min\_node\_size\_valor,}
      \AttributeTok{importance     =} \StringTok{"permutation"}\NormalTok{,}
      \AttributeTok{seed           =} \DecValTok{123}
\NormalTok{    )}

\NormalTok{    predicciones[i] }\OtherTok{\textless{}{-}} \FunctionTok{predict}\NormalTok{(modelo, }\AttributeTok{data =}\NormalTok{ prueba)}\SpecialCharTok{$}\NormalTok{predictions}
\NormalTok{  \}}

  \FunctionTok{list}\NormalTok{(}
    \AttributeTok{metricas =} \FunctionTok{calcular\_metricas}\NormalTok{(datos}\SpecialCharTok{$}\NormalTok{casos\_divorcio, predicciones) }\SpecialCharTok{\%\textgreater{}\%}
      \FunctionTok{mutate}\NormalTok{(}
        \AttributeTok{mtry          =}\NormalTok{ mtry\_valor,}
        \AttributeTok{min\_node\_size =}\NormalTok{ min\_node\_size\_valor,}
        \AttributeTok{num\_trees     =}\NormalTok{ num\_trees\_valor}
\NormalTok{      ),}
    \AttributeTok{predicciones =}\NormalTok{ predicciones}
\NormalTok{  )}
\NormalTok{\}}
\end{Highlighting}
\end{Shaded}

\begin{Shaded}
\begin{Highlighting}[]
\NormalTok{mtry\_base }\OtherTok{\textless{}{-}} \FunctionTok{max}\NormalTok{(}\DecValTok{1}\NormalTok{, }\FunctionTok{floor}\NormalTok{(}\FunctionTok{sqrt}\NormalTok{(}\FunctionTok{length}\NormalTok{(predictores\_rf))))}

\NormalTok{loo\_rf\_1 }\OtherTok{\textless{}{-}} \FunctionTok{loo\_rf}\NormalTok{(datos\_rf, }\AttributeTok{mtry\_valor =}\NormalTok{ mtry\_base,}
                   \AttributeTok{min\_node\_size\_valor =} \DecValTok{3}\NormalTok{, }\AttributeTok{num\_trees\_valor =} \DecValTok{300}\NormalTok{)}

\NormalTok{loo\_rf\_2 }\OtherTok{\textless{}{-}} \FunctionTok{loo\_rf}\NormalTok{(datos\_rf, }\AttributeTok{mtry\_valor =} \FunctionTok{max}\NormalTok{(}\DecValTok{1}\NormalTok{, mtry\_base }\SpecialCharTok{{-}} \DecValTok{1}\NormalTok{),}
                   \AttributeTok{min\_node\_size\_valor =} \DecValTok{4}\NormalTok{, }\AttributeTok{num\_trees\_valor =} \DecValTok{500}\NormalTok{)}

\NormalTok{loo\_rf\_3 }\OtherTok{\textless{}{-}} \FunctionTok{loo\_rf}\NormalTok{(datos\_rf, }\AttributeTok{mtry\_valor =} \FunctionTok{min}\NormalTok{(}\FunctionTok{length}\NormalTok{(predictores\_rf), mtry\_base }\SpecialCharTok{+} \DecValTok{1}\NormalTok{),}
                   \AttributeTok{min\_node\_size\_valor =} \DecValTok{2}\NormalTok{, }\AttributeTok{num\_trees\_valor =} \DecValTok{500}\NormalTok{)}

\NormalTok{pred\_loo\_rf }\OtherTok{\textless{}{-}} \FunctionTok{list}\NormalTok{(}
  \AttributeTok{RF\_1 =}\NormalTok{ loo\_rf\_1}\SpecialCharTok{$}\NormalTok{predicciones,}
  \AttributeTok{RF\_2 =}\NormalTok{ loo\_rf\_2}\SpecialCharTok{$}\NormalTok{predicciones,}
  \AttributeTok{RF\_3 =}\NormalTok{ loo\_rf\_3}\SpecialCharTok{$}\NormalTok{predicciones}
\NormalTok{)}

\NormalTok{metricas\_rf }\OtherTok{\textless{}{-}} \FunctionTok{bind\_rows}\NormalTok{(}
\NormalTok{  loo\_rf\_1}\SpecialCharTok{$}\NormalTok{metricas }\SpecialCharTok{\%\textgreater{}\%} \FunctionTok{mutate}\NormalTok{(}\AttributeTok{modelo =} \StringTok{"RF\_1"}\NormalTok{),}
\NormalTok{  loo\_rf\_2}\SpecialCharTok{$}\NormalTok{metricas }\SpecialCharTok{\%\textgreater{}\%} \FunctionTok{mutate}\NormalTok{(}\AttributeTok{modelo =} \StringTok{"RF\_2"}\NormalTok{),}
\NormalTok{  loo\_rf\_3}\SpecialCharTok{$}\NormalTok{metricas }\SpecialCharTok{\%\textgreater{}\%} \FunctionTok{mutate}\NormalTok{(}\AttributeTok{modelo =} \StringTok{"RF\_3"}\NormalTok{)}
\NormalTok{) }\SpecialCharTok{\%\textgreater{}\%}
  \FunctionTok{mutate}\NormalTok{(}
    \AttributeTok{algoritmo =} \StringTok{"Random Forest"}\NormalTok{,}
    \AttributeTok{rol       =} \StringTok{"Apoyo"}
\NormalTok{  ) }\SpecialCharTok{\%\textgreater{}\%}
  \FunctionTok{select}\NormalTok{(algoritmo, rol, modelo, mtry, min\_node\_size, num\_trees, MAE, RMSE, MAPE)}

\NormalTok{metricas\_rf }\SpecialCharTok{\%\textgreater{}\%}
  \FunctionTok{arrange}\NormalTok{(MAE, RMSE) }\SpecialCharTok{\%\textgreater{}\%}
  \FunctionTok{kable}\NormalTok{(}\AttributeTok{digits =} \DecValTok{3}\NormalTok{, }\AttributeTok{caption =} \StringTok{"Comparación de modelos Random Forest (LOOCV)"}\NormalTok{)}
\end{Highlighting}
\end{Shaded}

\begin{longtable}[]{@{}
  >{\raggedright\arraybackslash}p{(\columnwidth - 16\tabcolsep) * \real{0.1772}}
  >{\raggedright\arraybackslash}p{(\columnwidth - 16\tabcolsep) * \real{0.0759}}
  >{\raggedright\arraybackslash}p{(\columnwidth - 16\tabcolsep) * \real{0.0886}}
  >{\raggedleft\arraybackslash}p{(\columnwidth - 16\tabcolsep) * \real{0.0633}}
  >{\raggedleft\arraybackslash}p{(\columnwidth - 16\tabcolsep) * \real{0.1772}}
  >{\raggedleft\arraybackslash}p{(\columnwidth - 16\tabcolsep) * \real{0.1266}}
  >{\raggedleft\arraybackslash}p{(\columnwidth - 16\tabcolsep) * \real{0.1013}}
  >{\raggedleft\arraybackslash}p{(\columnwidth - 16\tabcolsep) * \real{0.1013}}
  >{\raggedleft\arraybackslash}p{(\columnwidth - 16\tabcolsep) * \real{0.0886}}@{}}
\caption{Comparación de modelos Random Forest (LOOCV)}\tabularnewline
\toprule\noalign{}
\begin{minipage}[b]{\linewidth}\raggedright
algoritmo
\end{minipage} & \begin{minipage}[b]{\linewidth}\raggedright
rol
\end{minipage} & \begin{minipage}[b]{\linewidth}\raggedright
modelo
\end{minipage} & \begin{minipage}[b]{\linewidth}\raggedleft
mtry
\end{minipage} & \begin{minipage}[b]{\linewidth}\raggedleft
min\_node\_size
\end{minipage} & \begin{minipage}[b]{\linewidth}\raggedleft
num\_trees
\end{minipage} & \begin{minipage}[b]{\linewidth}\raggedleft
MAE
\end{minipage} & \begin{minipage}[b]{\linewidth}\raggedleft
RMSE
\end{minipage} & \begin{minipage}[b]{\linewidth}\raggedleft
MAPE
\end{minipage} \\
\midrule\noalign{}
\endfirsthead
\toprule\noalign{}
\begin{minipage}[b]{\linewidth}\raggedright
algoritmo
\end{minipage} & \begin{minipage}[b]{\linewidth}\raggedright
rol
\end{minipage} & \begin{minipage}[b]{\linewidth}\raggedright
modelo
\end{minipage} & \begin{minipage}[b]{\linewidth}\raggedleft
mtry
\end{minipage} & \begin{minipage}[b]{\linewidth}\raggedleft
min\_node\_size
\end{minipage} & \begin{minipage}[b]{\linewidth}\raggedleft
num\_trees
\end{minipage} & \begin{minipage}[b]{\linewidth}\raggedleft
MAE
\end{minipage} & \begin{minipage}[b]{\linewidth}\raggedleft
RMSE
\end{minipage} & \begin{minipage}[b]{\linewidth}\raggedleft
MAPE
\end{minipage} \\
\midrule\noalign{}
\endhead
\bottomrule\noalign{}
\endlastfoot
Random Forest & Apoyo & RF\_3 & 3 & 2 & 500 & 300.120 & 661.341 &
75.201 \\
Random Forest & Apoyo & RF\_2 & 1 & 4 & 500 & 301.877 & 668.170 &
75.212 \\
Random Forest & Apoyo & RF\_1 & 2 & 3 & 300 & 307.793 & 665.297 &
77.588 \\
\end{longtable}

La tabla resume tres configuraciones que difieren en \texttt{mtry},
\texttt{min.node.size} y número de árboles. La variante con menor
\textbf{MAE} y \textbf{RMSE} bajo LOOCV se adopta como Random Forest
final; si dos filas empatan, el orden \texttt{arrange(MAE,\ RMSE)}
prioriza el error absoluto medio.

\begin{Shaded}
\begin{Highlighting}[]
\NormalTok{mejor\_rf }\OtherTok{\textless{}{-}}\NormalTok{ metricas\_rf }\SpecialCharTok{\%\textgreater{}\%}
  \FunctionTok{arrange}\NormalTok{(MAE, RMSE) }\SpecialCharTok{\%\textgreater{}\%}
  \FunctionTok{slice}\NormalTok{(}\DecValTok{1}\NormalTok{)}

\NormalTok{modelo\_rf\_final }\OtherTok{\textless{}{-}} \FunctionTok{ranger}\NormalTok{(}
\NormalTok{  casos\_divorcio }\SpecialCharTok{\textasciitilde{}}\NormalTok{ .,}
  \AttributeTok{data          =}\NormalTok{ datos\_rf,}
  \AttributeTok{num.trees     =}\NormalTok{ mejor\_rf}\SpecialCharTok{$}\NormalTok{num\_trees,}
  \AttributeTok{mtry          =}\NormalTok{ mejor\_rf}\SpecialCharTok{$}\NormalTok{mtry,}
  \AttributeTok{min.node.size =}\NormalTok{ mejor\_rf}\SpecialCharTok{$}\NormalTok{min\_node\_size,}
  \AttributeTok{importance    =} \StringTok{"permutation"}\NormalTok{,}
  \AttributeTok{seed          =} \DecValTok{123}
\NormalTok{)}

\NormalTok{pred\_rf\_final     }\OtherTok{\textless{}{-}} \FunctionTok{predict}\NormalTok{(modelo\_rf\_final, }\AttributeTok{data =}\NormalTok{ datos\_rf)}\SpecialCharTok{$}\NormalTok{predictions}
\NormalTok{metricas\_rf\_final }\OtherTok{\textless{}{-}} \FunctionTok{calcular\_metricas}\NormalTok{(datos\_rf}\SpecialCharTok{$}\NormalTok{casos\_divorcio, pred\_rf\_final)}
\end{Highlighting}
\end{Shaded}

\begin{Shaded}
\begin{Highlighting}[]
\NormalTok{diagnostico\_rf }\OtherTok{\textless{}{-}}\NormalTok{ datos\_modelo }\SpecialCharTok{\%\textgreater{}\%}
  \FunctionTok{select}\NormalTok{(cod\_depto, casos\_divorcio) }\SpecialCharTok{\%\textgreater{}\%}
  \FunctionTok{mutate}\NormalTok{(}\AttributeTok{pred\_rf =}\NormalTok{ pred\_rf\_final)}

\NormalTok{diagnostico\_rf }\SpecialCharTok{\%\textgreater{}\%}
  \FunctionTok{kable}\NormalTok{(}\AttributeTok{digits =} \DecValTok{3}\NormalTok{, }\AttributeTok{caption =} \StringTok{"Predicciones del modelo Random Forest final por departamento"}\NormalTok{)}
\end{Highlighting}
\end{Shaded}

\begin{longtable}[]{@{}rrr@{}}
\caption{Predicciones del modelo Random Forest final por
departamento}\tabularnewline
\toprule\noalign{}
cod\_depto & casos\_divorcio & pred\_rf \\
\midrule\noalign{}
\endfirsthead
\toprule\noalign{}
cod\_depto & casos\_divorcio & pred\_rf \\
\midrule\noalign{}
\endhead
\bottomrule\noalign{}
\endlastfoot
1 & 3295 & 2125.349 \\
2 & 179 & 228.292 \\
3 & 174 & 304.661 \\
4 & 210 & 323.432 \\
5 & 497 & 459.001 \\
6 & 316 & 328.707 \\
7 & 137 & 360.830 \\
8 & 201 & 250.361 \\
9 & 794 & 573.675 \\
10 & 485 & 459.975 \\
11 & 331 & 339.350 \\
12 & 565 & 476.578 \\
13 & 450 & 394.351 \\
14 & 314 & 313.666 \\
15 & 145 & 204.933 \\
16 & 198 & 264.367 \\
17 & 326 & 378.837 \\
18 & 286 & 285.284 \\
19 & 195 & 206.618 \\
20 & 212 & 233.076 \\
21 & 238 & 247.769 \\
22 & 402 & 340.639 \\
\end{longtable}

La tabla alinea \texttt{cod\_depto} desde \texttt{datos\_modelo} con las
predicciones del bosque; el orden de filas coincide con
\texttt{datos\_rf} porque ambos provienen del mismo
\texttt{datos\_modelo} sin reordenar.

\begin{Shaded}
\begin{Highlighting}[]
\FunctionTok{ggplot}\NormalTok{(diagnostico\_rf, }\FunctionTok{aes}\NormalTok{(}\AttributeTok{x =}\NormalTok{ pred\_rf, }\AttributeTok{y =}\NormalTok{ casos\_divorcio)) }\SpecialCharTok{+}
  \FunctionTok{geom\_point}\NormalTok{(}\AttributeTok{size =} \FloatTok{2.5}\NormalTok{) }\SpecialCharTok{+}
  \FunctionTok{geom\_abline}\NormalTok{(}\AttributeTok{slope =} \DecValTok{1}\NormalTok{, }\AttributeTok{intercept =} \DecValTok{0}\NormalTok{, }\AttributeTok{linetype =} \StringTok{"dashed"}\NormalTok{, }\AttributeTok{color =} \StringTok{"gray50"}\NormalTok{) }\SpecialCharTok{+}
  \FunctionTok{labs}\NormalTok{(}
    \AttributeTok{title =} \StringTok{"Valores observados vs predichos — Random Forest"}\NormalTok{,}
    \AttributeTok{x     =} \StringTok{"Casos predichos"}\NormalTok{,}
    \AttributeTok{y     =} \StringTok{"Casos observados"}
\NormalTok{  ) }\SpecialCharTok{+}
  \FunctionTok{theme\_minimal}\NormalTok{()}
\end{Highlighting}
\end{Shaded}

\begin{figure}
\centering
\includegraphics{Proyecto_Entrega2_Presentacion_Resultados_files/figure-latex/rf-grafica-obs-pred-1.pdf}
\caption{Valores observados frente a predichos para Random Forest
(ajuste en muestra completa). La línea discontinua es predicción
perfecta.}
\end{figure}

El gráfico usa la misma escala que NB y Poisson (divorcios observados vs
predichos). Random Forest puede ajustar mejor algunos puntos del
interior si existen interacciones no modeladas por los GLM, pero con
\(n\) tan pequeño las diferencias deben interpretarse con cautela.

\begin{Shaded}
\begin{Highlighting}[]
\NormalTok{importancia\_rf }\OtherTok{\textless{}{-}} \FunctionTok{tibble}\NormalTok{(}
  \AttributeTok{variable    =} \FunctionTok{names}\NormalTok{(modelo\_rf\_final}\SpecialCharTok{$}\NormalTok{variable.importance),}
  \AttributeTok{importancia =} \FunctionTok{as.numeric}\NormalTok{(modelo\_rf\_final}\SpecialCharTok{$}\NormalTok{variable.importance)}
\NormalTok{) }\SpecialCharTok{\%\textgreater{}\%}
  \FunctionTok{arrange}\NormalTok{(}\FunctionTok{desc}\NormalTok{(importancia))}

\NormalTok{importancia\_rf }\SpecialCharTok{\%\textgreater{}\%}
  \FunctionTok{kable}\NormalTok{(}\AttributeTok{digits =} \DecValTok{3}\NormalTok{, }\AttributeTok{caption =} \StringTok{"Importancia de variables en Random Forest (permutación)"}\NormalTok{)}
\end{Highlighting}
\end{Shaded}

\begin{longtable}[]{@{}lr@{}}
\caption{Importancia de variables en Random Forest
(permutación)}\tabularnewline
\toprule\noalign{}
variable & importancia \\
\midrule\noalign{}
\endfirsthead
\toprule\noalign{}
variable & importancia \\
\midrule\noalign{}
\endhead
\bottomrule\noalign{}
\endlastfoot
casos\_violencia\_z & 13975.365 \\
es\_capital & 0.000 \\
edad\_victima\_prom\_z & -4107.243 \\
porc\_urbano\_z & -4310.927 \\
\end{longtable}

\begin{Shaded}
\begin{Highlighting}[]
\FunctionTok{ggplot}\NormalTok{(importancia\_rf, }\FunctionTok{aes}\NormalTok{(}\AttributeTok{x =} \FunctionTok{reorder}\NormalTok{(variable, importancia), }\AttributeTok{y =}\NormalTok{ importancia)) }\SpecialCharTok{+}
  \FunctionTok{geom\_col}\NormalTok{(}\AttributeTok{fill =} \StringTok{"steelblue"}\NormalTok{, }\AttributeTok{alpha =} \FloatTok{0.85}\NormalTok{) }\SpecialCharTok{+}
  \FunctionTok{coord\_flip}\NormalTok{() }\SpecialCharTok{+}
  \FunctionTok{labs}\NormalTok{(}
    \AttributeTok{title =} \StringTok{"Importancia de variables — Random Forest"}\NormalTok{,}
    \AttributeTok{x     =} \StringTok{"Variable"}\NormalTok{,}
    \AttributeTok{y     =} \StringTok{"Importancia por permutación"}
\NormalTok{  ) }\SpecialCharTok{+}
  \FunctionTok{theme\_minimal}\NormalTok{()}
\end{Highlighting}
\end{Shaded}

\begin{figure}
\centering
\includegraphics{Proyecto_Entrega2_Presentacion_Resultados_files/figure-latex/rf-importancia-grafica-1.pdf}
\caption{Importancia por permutación en el Random Forest final (mayor
valor = más contribución al error cuando la variable se permuta).}
\end{figure}

Tras ejecutar el documento, la variable con mayor importancia de
permutación indica qué predictor el bosque usa con más intensidad para
reducir el error de predicción. Si \textbf{\texttt{es\_capital}} y/o
\textbf{\texttt{casos\_violencia\_z}} lideran la importancia, el
resultado es coherente con la Regresión Binomial Negativa, donde la
dummy metropolitana y el volumen de violencia concentran buena parte de
la explicación a nivel departamental. Si otra variable gana claramente,
sugiere que Random Forest está captando estructura adicional (por
ejemplo, combinaciones no lineales con urbanización o edad) que el GLM
lineal en los predictores no incorpora; ello no invalida el modelo
principal, pero enriquece la lectura exploratoria.

Dado el tamaño muestral, Random Forest tiene \textbf{mayor riesgo de
ajuste específico a estos 22 puntos} que la Binomial Negativa; por eso
se mantiene como apoyo y no como modelo de decisión principal.

\begin{Shaded}
\begin{Highlighting}[]
\FunctionTok{stopifnot}\NormalTok{(}\FunctionTok{exists}\NormalTok{(}\StringTok{"metricas\_rf"}\NormalTok{))}
\FunctionTok{stopifnot}\NormalTok{(}\FunctionTok{exists}\NormalTok{(}\StringTok{"modelo\_rf\_final"}\NormalTok{))}
\FunctionTok{stopifnot}\NormalTok{(}\FunctionTok{exists}\NormalTok{(}\StringTok{"pred\_rf\_final"}\NormalTok{))}
\FunctionTok{stopifnot}\NormalTok{(}\FunctionTok{exists}\NormalTok{(}\StringTok{"diagnostico\_rf"}\NormalTok{))}
\FunctionTok{stopifnot}\NormalTok{(}\FunctionTok{nrow}\NormalTok{(metricas\_rf) }\SpecialCharTok{\textgreater{}=} \DecValTok{3}\NormalTok{)}
\FunctionTok{stopifnot}\NormalTok{(}\FunctionTok{nrow}\NormalTok{(datos\_rf) }\SpecialCharTok{==} \FunctionTok{nrow}\NormalTok{(datos\_modelo))}
\FunctionTok{cat}\NormalTok{(}\StringTok{"Random Forest (apoyo) OK — configuración final: mtry ="}\NormalTok{, mejor\_rf}\SpecialCharTok{$}\NormalTok{mtry,}
    \StringTok{", min.node.size ="}\NormalTok{, mejor\_rf}\SpecialCharTok{$}\NormalTok{min\_node\_size,}
    \StringTok{", num.trees ="}\NormalTok{, mejor\_rf}\SpecialCharTok{$}\NormalTok{num\_trees, }\StringTok{"}\SpecialCharTok{\textbackslash{}n}\StringTok{"}\NormalTok{)}
\end{Highlighting}
\end{Shaded}

\begin{verbatim}
## Random Forest (apoyo) OK — configuración final: mtry = 3 , min.node.size = 2 , num.trees = 500
\end{verbatim}

\end{document}
